
% ============================================
% chapter12_full_zh.tex — 第12章：主理想整环上的模
% 按原书第12章 PDF 逐节重译校订；正文、证明、例子与练习均保留
% ============================================

\chapter{主理想整环上的模}

本章的主要目的是证明有限生成模在一类特别好的环上，即在主理想整环（Principal Ideal Domain, P.I.D.）上的结构定理。这个定理是环的理想结构（对 P.I.D. 来说特别简单）反映在它的模结构中的一个例子。如果把这个结果应用于 P.I.D. 为整数环 $\mathbb Z$ 的情形，那么就得到有限生成阿贝尔群基本定理的一个证明（我们在第 5 章已经不加证明地考察过它）。如果把这个结构定理应用于 P.I.D. 为域 $F$ 上关于 $x$ 的多项式环 $F[x]$ 的情形，那么就得到所谓矩阵有理标准形与 Jordan 标准形的基本结果。在进入证明以前，我们先简要讨论这两个重要应用。

我们在第 5 章已经讨论过如下结果：任意有限生成阿贝尔群同构于若干循环阿贝尔群的直和，而这些循环阿贝尔群要么是 $\mathbb Z$，要么是某个正整数 $n\neq 0$ 对应的 $\mathbb Z/n\mathbb Z$。还要回忆，阿贝尔群与 $\mathbb Z$-模是同一件事。由于 $\mathbb Z$ 的理想恰好是平凡理想 $(0)$ 以及由正整数 $n$ 生成的主理想 $(n)=n\mathbb Z$，我们看到，用模的语言来说，有限生成阿贝尔群基本定理断言任意有限生成 $\mathbb Z$-模都是形如 $\mathbb Z/I$ 的模（其中 $I$ 是 $\mathbb Z$ 的理想；这些正是循环 $\mathbb Z$-模）的直和，并且当直和写成一种特定形式时还带有唯一性陈述。注意 $\mathbb Z$ 的理想结构与它的有限生成模（也就是有限生成阿贝尔群）的结构之间的对应关系。

P.I.D. 上有限生成模基本定理说，当主理想整环 $\mathbb Z$ 被任意 P.I.D. 取代时，同样的结果仍成立。特别地，我们已经在第 10 章看到，关于域 $F$ 上多项式环 $F[x]$ 的模与一个 $F$-向量空间 $V$ 加上一个固定的线性变换 $T$ 是同一件事（其中元素 $x$ 在 $V$ 上按线性变换 $T$ 作用）。在这个情形中，基本定理将说明这样的向量空间是形如 $F[x]/I$ 的模的直和，其中 $I$ 是 $F[x]$ 的理想，因而或者是平凡理想 $(0)$，或者是由某个非零多项式 $f(x)$ 生成的主理想 $(f(x))$（这些就是循环 $F[x]$-模）；同样地，当直和写成一种特定形式时还带有唯一性陈述。如果把这翻译回向量空间和线性变换的语言，我们便能得到关于线性变换 $T$ 的信息。

例如，设 $V$ 是 $F$ 上 $n$ 维向量空间并且我们为 $V$ 选定一组基。那么给出 $V$ 到自身的线性变换 $T$，等价于给出一个系数在 $F$ 中的 $n\times n$ 矩阵 $A$（如果为 $V$ 选取另一组基，则得到 $T$ 的另一个矩阵 $B$，它与 $A$ 相似，即 $B$ 具有 $P^{-1}AP$ 的形式，其中 $P$ 是定义换基的可逆矩阵）。我们将看到，在这种情形下，基本定理蕴含（假设域 $F$ 含有给定线性变换 $T$ 的所有“特征值”）存在 $V$ 的一组基，使得 $T$ 的相应矩阵尽可能接近对角矩阵，因而有一个特别简单的形式。这就是 Jordan 标准形。有理标准形是 $T$ 的矩阵的另一种简单形式（它不要求 $T$ 的特征值属于 $F$）。这样，我们就能给出域 $F$ 上任意 $n\times n$ 矩阵的标准形，即找到与给定 $n\times n$ 矩阵相似而又特别简单的矩阵（例如几乎对角）。

\paragraph{例}
令
\[
V=\mathbb Q^3=\{(x,y,z)\mid x,y,z\in\mathbb Q\}
\]
为有理数域 $F=\mathbb Q$ 上通常的三维有序三元组向量空间，并设 $T$ 是线性变换
\[
T(x,y,z)=(9x+4y+5z,-4x-3z,-6x-4y-2z),\qquad x,y,z\in\mathbb Q.
\]
如果取 $V$ 的标准基 $e_1=(1,0,0)$，$e_2=(0,1,0)$，$e_3=(0,0,1)$，则表示这个线性变换的矩阵为
\[
A=\begin{pmatrix}
9&4&5\\
-4&0&-3\\
-6&-4&-2
\end{pmatrix}.
\]
我们将看到，这个矩阵 $A$ 的 Jordan 标准形是更简单的矩阵
\[
B=\begin{pmatrix}
2&1&0\\
0&2&0\\
0&0&3
\end{pmatrix},
\]
它通过改取基 $f_1=(2,-1,-2)$，$f_2=(1,0,-1)$，$f_3=(3,-2,-2)$ 得到，因为在这种情况下
\[
\begin{aligned}
T(f_1)&=T(2,-1,-2)=(4,-2,-4)=2f_1+0f_2+0f_3,\\
T(f_2)&=T(1,0,-1)=(4,-1,-4)=1f_1+2f_2+0f_3,\\
T(f_3)&=T(3,-2,-2)=(9,-6,-6)=0f_1+0f_2+3f_3.
\end{aligned}
\]
所以相对于这组基表示 $T$ 的矩阵的列分别是 $(2,0,0)$、$(1,2,0)$ 与 $(0,0,3)$，也就是说 $T$ 相对于这组基的矩阵是 $B$。特别地，$A$ 与更简单的矩阵 $B$ 相似。事实上这个线性变换 $T$ 不能对角化（即不存在一组 $V$ 的基使相应矩阵为对角矩阵），所以矩阵 $B$ 已经是 $T$ 可能达到的最接近对角矩阵的形式。

下面第一节给出若干一般定义，并陈述和证明任意 P.I.D. 上的基本定理；然后再回到标准形的应用（对阿贝尔群的应用出现在第 5 章）。这些应用可以独立于一般证明来阅读。对 Euclidean 整环有效的一种替代证明（因而特别适用于环 $\mathbb Z$ 和 $F[x]$），按行、列运算的方式给出，并且在计算上有用；它在练习中加以概述。

\section{基本理论}

我们首先描述若干一般的有限性条件。令 $R$ 为环，$M$ 为左 $R$-模。

\noindent\textbf{定义.}
\begin{enumerate}[label=(\arabic*)]
\item 左 $R$-模 $M$ 称为 \emph{Noether $R$-模}，或称满足子模上的升链条件（ascending chain condition on submodules，简称 A.C.C. on submodules），如果不存在子模的无限严格递增链；也就是说，每当
\[
M_1\subseteq M_2\subseteq M_3\subseteq\cdots
\]
是 $M$ 的子模升链时，存在正整数 $m$，使得对所有 $k\geq m$ 都有 $M_k=M_m$（因此这条链在第 $m$ 阶段稳定：$M_m=M_{m+1}=M_{m+2}=\cdots$）。
\item 环 $R$ 称为 \emph{Noether 环}，如果 $R$ 作为它自身上的左模是 Noether 的；也就是说，如果 $R$ 中不存在左理想的无限严格递增链。
\end{enumerate}

在一个（可能非交换的）环 $R$ 中，也可以类似地表述右理想和双边理想上的 A.C.C. 概念。对非交换环而言，这些性质不必相互相关。

\noindent\textbf{定理 1.} 设 $R$ 为环，$M$ 为左 $R$-模。下列条件等价：
\begin{enumerate}[label=(\arabic*)]
\item $M$ 是 Noether $R$-模；
\item $M$ 的每个非空子模集合在包含关系下含有极大元；
\item $M$ 的每个子模都是有限生成的。
\end{enumerate}

\begin{proof}
$[(1)\Rightarrow(2)]$ 假设 $M$ 是 Noether 的，并令 $\Sigma$ 是 $M$ 的任意非空子模族。任选 $M_1\in\Sigma$。如果 $M_1$ 已是 $\Sigma$ 中的极大元，则 (2) 成立；于是设 $M_1$ 不是极大元。则存在某个 $M_2\in\Sigma$ 使 $M_1\subsetneq M_2$。如果 $M_2$ 是 $\Sigma$ 中的极大元，则 (2) 成立；否则又存在 $M_3\in\Sigma$ 严格包含 $M_2$。如此继续下去可见，若 (2) 失败，则可借助选择公理构造出 $\Sigma$ 中元素的无限严格递增链
\[
M_1\subsetneq M_2\subsetneq M_3\subsetneq\cdots,
\]
这与 (1) 矛盾。

$[(2)\Rightarrow(3)]$ 假设 (2) 成立，令 $N$ 是 $M$ 的任意子模。令 $\Sigma$ 为 $N$ 的所有有限生成子模的集合。由于 $\{0\}\in\Sigma$，这个集合非空。由 (2)，$\Sigma$ 含有极大元 $N'$。若 $N'\neq N$，取 $x\in N-N'$。由于 $N'\in\Sigma$，子模 $N'$ 按假设有限生成，因此由 $N'$ 与 $x$ 生成的子模也有限生成，并且严格包含 $N'$。这与 $N'$ 的极大性矛盾。因此 $N=N'$ 有限生成。

$[(3)\Rightarrow(1)]$ 假设 (3) 成立，并令
\[
M_1\subseteq M_2\subseteq M_3\subseteq\cdots
\]
是 $M$ 的子模链。令
\[
N=\bigcup_{i=1}^{\infty}M_i,
\]
注意 $N$ 是一个子模。由 (3)，$N$ 有有限个生成元，设为 $a_1,a_2,\ldots,a_n$。因为每个 $a_i\in N$，每个 $a_i$ 都属于这条链中的某个子模，设 $a_i\in M_{j_i}$。令
\[
m=\max\{j_1,j_2,\ldots,j_n\}.
\]
则所有 $a_i$ 都属于 $M_m$，所以它们生成的模包含在 $M_m$ 中，即 $N\subseteq M_m$。于是 $M_m=N=M_k$ 对所有 $k\geq m$ 成立，这证明了 (1)。
\end{proof}

\noindent\textbf{推论 2.} 若 $R$ 是 P.I.D.，则 $R$ 的每个非空理想集合有极大元，并且 $R$ 是 Noether 环。

\begin{proof}
P.I.D. $R$ 在定理 1 中对 $M=R$ 满足条件 (3)，因为 $R$ 的每个理想都是主理想，因而由一个元素生成。
\end{proof}

回忆：即使 $M$ 本身是有限生成 $R$-模，$M$ 的子模也未必有限生成，所以 $M$ 是 Noether $R$-模这个条件通常强于 $M$ 是有限生成 $R$-模这个条件。

在转向本章主要结果之前，我们需要一个关于“线性相关”的结果。

\noindent\textbf{命题 3.} 设 $R$ 是整环，$M$ 是秩 $n<\infty$ 的自由 $R$-模。则 $M$ 中任意 $n+1$ 个元素都是 $R$-线性相关的；也就是说，对任意 $y_1,y_2,\ldots,y_{n+1}\in M$，存在不全为零的 $r_1,r_2,\ldots,r_{n+1}\in R$，使得
\[
r_1y_1+r_2y_2+\cdots+r_{n+1}y_{n+1}=0.
\]

\begin{proof}
最快的证明是把 $R$ 嵌入它的分式域 $F$（因为 $R$ 是整环），并观察到：由于
\[
M\cong R\oplus R\oplus\cdots\oplus R\quad(n\text{ 次}),
\]
我们得到
\[
F\otimes_R M\cong F\oplus F\oplus\cdots\oplus F.
\]
后者是 $F$ 上的 $n$ 维向量空间，所以 $M$ 中任意 $n+1$ 个元素在张量到 $F$ 后都是 $F$-线性相关的。把这些标量的分母清除（例如同乘所有分母的乘积）便得到这 $n+1$ 个元素之间的一个 $R$-线性相关关系。

也可如下证明。令 $e_1,\ldots,e_n$ 是自由 $R$-模 $M$ 的一组基，并令 $y_1,\ldots,y_{n+1}$ 是 $M$ 的任意 $n+1$ 个元素。对 $1\leq i\leq n+1$，把
\[
y_i=a_{1i}e_1+a_{2i}e_2+\cdots+a_{ni}e_n
\]
写在基 $e_1,e_2,\ldots,e_n$ 下。令 $A$ 为 $(n+1)\times(n+1)$ 矩阵，它的第 $i,j$ 个条目为 $a_{ij}$（$1\leq i\leq n$，$1\leq j\leq n+1$），最后一行全为 $0$，所以当然 $\det A=0$。由于 $R$ 是整环，第 11.4 节推论 27 说明 $A$ 的列 $R$-线性相关。列之间的任一线性相关关系正给出 $y_i$ 之间的相应线性相关关系，从而完成证明。
\end{proof}

若 $R$ 是任意整环，$M$ 是任意 $R$-模，回忆
\[
\operatorname{Tor}(M)=\{x\in M\mid rx=0\text{ 对某个非零 }r\in R\}
\]
是 $M$ 的子模（称为 $M$ 的\emph{挠子模}）。如果 $N$ 是 $\operatorname{Tor}(M)$ 的任意子模，则称 $N$ 为 $M$ 的一个挠子模（因此 $\operatorname{Tor}(M)$ 是 $M$ 的所有挠子模的并，即 $M$ 的极大挠子模）。如果 $\operatorname{Tor}(M)=0$，则称 $M$ 是\emph{无挠}的。

对 $M$ 的任意子模 $N$，$N$ 的零化子是 $R$ 的理想
\[
\operatorname{Ann}(N)=\{r\in R\mid rn=0\text{ 对所有 }n\in N\}.
\]
注意，如果 $N$ 不是 $M$ 的挠子模，则 $\operatorname{Ann}(N)=(0)$。容易看出，如果 $N,L$ 是 $M$ 的子模且 $N\subseteq L$，则 $\operatorname{Ann}(L)\subseteq \operatorname{Ann}(N)$。若 $R$ 是 P.I.D. 且 $N\subseteq L\subseteq M$，并且 $\operatorname{Ann}(N)=(a)$、$\operatorname{Ann}(L)=(b)$，则 $a\mid b$。特别地，$M$ 的任一元素 $x$ 的零化子整除 $M$ 的零化子（当 $R=\mathbb Z$ 时，这由 Lagrange 定理蕴含）。

\noindent\textbf{定义.} 对任意整环 $R$，$R$-模 $M$ 的\emph{秩}是 $M$ 中 $R$-线性无关元素个数的最大值。

前一命题说明：在整环上，一个有限秩自由 $R$-模 $M$ 的子模的秩被 $M$ 的秩所界定。这个秩的概念与此前同名术语一致。如果环 $R=F$ 是域，则 $R$-模 $M$ 的秩就是 $M$ 作为 $F$-向量空间的维数，并且任意极大 $F$-线性无关集都是 $M$ 的一组基。然而，对一般整环，秩为 $n$ 的 $R$-模 $M$ 不必有“基”，也就是说，即使 $M$ 无挠，它也不必是自由 $R$-模。因此使用秩的概念时需要谨慎，尤其要注意 $M$ 的挠元素。练习 1--6 与 20 给出秩的另一种刻画，并给出一些（无挠）秩 $1$ 但非自由的 $R$-模例子。

下面这个重要结果说明：如果 $N$ 是 P.I.D. 上有限秩自由模的子模，则 $N$ 仍是有限秩自由模；而且可以为两个模选择生成元，使它们之间具有一种简单关系。

\noindent\textbf{定理 4.} 设 $R$ 是主理想整环，$M$ 是有限秩 $n$ 的自由 $R$-模，$N$ 是 $M$ 的子模。则
\begin{enumerate}[label=(\arabic*)]
\item $N$ 是自由模，秩为 $m$，且 $m\leq n$；
\item 存在 $M$ 的一组基 $y_1,y_2,\ldots,y_n$，使得 $a_1y_1,a_2y_2,\ldots,a_my_m$ 是 $N$ 的一组基，其中 $a_1,a_2,\ldots,a_m$ 是 $R$ 中非零元素，并满足整除关系
\[
a_1\mid a_2\mid\cdots\mid a_m.
\]
\end{enumerate}

\begin{proof}
当 $N=\{0\}$ 时定理平凡，故设 $N\neq\{0\}$。对每个 $R$-模同态 $\varphi:M\to R$，$N$ 的像 $\varphi(N)$ 是 $R$ 的子模，即 $R$ 的理想。由于 $R$ 是 P.I.D.，该理想必为主理想，记为 $\varphi(N)=(a_\varphi)$，其中 $a_\varphi\in R$。令
\[
\Sigma=\{(a_\varphi)\mid \varphi\in\operatorname{Hom}_R(M,R)\}
\]
为以这种方式由 $M$ 到 $R$ 的 $R$-模同态得到的主理想集合。集合 $\Sigma$ 非空，因为取零同态便得到 $(0)\in\Sigma$。由推论 2，$\Sigma$ 至少有一个极大元；也就是说，存在至少一个同态 $v:M\to R$，使得主理想 $v(N)=(a_v)$ 不真包含在 $\Sigma$ 的任何其他元素中。令 $a_1=a_v$，并取 $y\in N$ 使 $v(y)=a_1$。

先说明 $a_1$ 非零。令 $x_1,x_2,\ldots,x_n$ 是自由模 $M$ 的任意一组基，并令 $\pi_i\in\operatorname{Hom}_R(M,R)$ 为相对于这组基投到第 $i$ 个坐标的自然投影同态。因为 $N\neq\{0\}$，存在某个 $i$ 使 $\pi_i(N)\neq0$；这特别说明 $\Sigma$ 不仅仅含有零理想。由于 $(a_1)$ 是 $\Sigma$ 的极大元，得到 $a_1\neq0$。

接下来说明，对每个 $\varphi\in\operatorname{Hom}_R(M,R)$，元素 $a_1$ 都整除 $\varphi(y)$。令 $d$ 是由 $a_1$ 和 $\varphi(y)$ 生成的主理想的生成元。因此 $d$ 是 $a_1$ 与 $\varphi(y)$ 的公因子，并且存在 $r_1,r_2\in R$ 使
\[
d=r_1a_1+r_2\varphi(y).
\]
考虑同态 $\psi=r_1v+r_2\varphi:M\to R$。则
\[
\psi(y)=(r_1v+r_2\varphi)(y)=r_1a_1+r_2\varphi(y)=d.
\]
所以 $d\in\psi(N)$，从而 $(d)\subseteq\psi(N)$。另一方面，因为 $d$ 是 $a_1$ 的因子，也有 $(a_1)\subseteq(d)$。于是
\[
(a_1)\subseteq(d)\subseteq\psi(N).
\]
由 $(a_1)$ 在 $\Sigma$ 中的极大性，只能有等号：$(a_1)=(d)=\psi(N)$。特别地，$(a_1)=(d)$ 说明 $a_1\mid\varphi(y)$，因为 $d$ 整除 $\varphi(y)$。

把这个结论应用于投影同态 $\pi_i$，可知 $a_1$ 整除所有 $\pi_i(y)$。写
\[
\pi_i(y)=a_1b_i\qquad(1\leq i\leq n)
\]
并定义
\[
y_1=\sum_{i=1}^n b_ix_i.
\]
注意 $a_1y_1=y$。因为
\[
a_1=v(y)=v(a_1y_1)=a_1v(y_1),
\]
且 $a_1$ 是整环 $R$ 中的非零元素，所以
\[
v(y_1)=1.
\]

现在验证这个元素 $y_1$ 可以取为 $M$ 的一组基中的一个元素，而且 $a_1y_1$ 可以取为 $N$ 的一组基中的一个元素；也就是证明
\[
\text{(a)}\quad M=Ry_1\oplus\ker v,
\qquad
\text{(b)}\quad N=Ra_1y_1\oplus (N\cap\ker v).
\]
为证明 (a)，取任意 $x\in M$，写成
\[
x=v(x)y_1+(x-v(x)y_1).
\]
因为
\[
v(x-v(x)y_1)=v(x)-v(x)v(y_1)=v(x)-v(x)\cdot1=0,
\]
可知 $x-v(x)y_1\in\ker v$。这说明 $x$ 可以写成 $Ry_1$ 中元素与 $\ker v$ 中元素之和，故 $M=Ry_1+\ker v$。若 $ry_1\in\ker v$，则
\[
0=v(ry_1)=rv(y_1)=r,
\]
所以这个交中的元素只能为 $0$，于是直和成立。

为证明 (b)，注意对每个 $x'\in N$，按 $a_1$ 作为 $v(N)$ 的生成元的定义，$v(x')$ 被 $a_1$ 整除。写 $v(x')=ba_1$（$b\in R$），则上面 (a) 中使用的分解给出
\[
x'=v(x')y_1+(x'-v(x')y_1)=ba_1y_1+(x'-ba_1y_1),
\]
其中第二项既在 $\ker v$ 中，也在 $N$ 中。这说明 $N=Ra_1y_1+(N\cap\ker v)$。该和为直和则是 (a) 中直和性的特例。

现在对 $N$ 的秩 $m$ 作归纳来证明 (1)。如果 $m=0$，则 $N$ 是挠模；但自由模无挠，故 $N=0$，(1) 平凡成立。于是设 $m>0$。由 (b) 中的直和容易看出 $N\cap\ker v$ 的秩为 $m-1$（参见练习 3）。由归纳假设，$N\cap\ker v$ 是秩 $m-1$ 的自由 $R$-模。再由 (b) 的直和性可知，把 $a_1y_1$ 添加到 $N\cap\ker v$ 的任一组基中，得到 $N$ 的一组基，所以 $N$ 也是自由的（秩为 $m$），这证明了 (1)。

最后对 $M$ 的秩 $n$ 作归纳证明 (2)。把 (1) 应用于子模 $\ker v$，可知这个子模是自由的；又由 (a) 的直和可知它的秩为 $n-1$。由归纳假设应用于模 $\ker v$（它扮演 $M$ 的角色）及其子模 $\ker v\cap N$（它扮演 $N$ 的角色），可知存在 $\ker v$ 的一组基 $y_2,y_3,\ldots,y_n$，使得 $a_2y_2,a_3y_3,\ldots,a_my_m$ 是 $N\cap\ker v$ 的一组基，并且 $a_2,a_3,\ldots,a_m\in R$ 满足
\[
a_2\mid a_3\mid\cdots\mid a_m.
\]
由 (a) 与 (b) 的直和性，$y_1,y_2,\ldots,y_n$ 是 $M$ 的一组基，而 $a_1y_1,a_2y_2,\ldots,a_my_m$ 是 $N$ 的一组基。为完成归纳，只需说明 $a_1\mid a_2$。定义同态 $\varphi:M\to R$，在 $M$ 的这组基上令
\[
\varphi(y_1)=\varphi(y_2)=1,
\qquad
\varphi(y_i)=0\quad(i>2).
\]
于是对这个同态 $\varphi$，有 $a_1=\varphi(a_1y_1)$，所以 $a_1\in\varphi(N)$，从而 $(a_1)\subseteq\varphi(N)$。由 $(a_1)$ 在 $\Sigma$ 中的极大性，得到 $(a_1)=\varphi(N)$。又因为 $a_2=\varphi(a_2y_2)\in\varphi(N)$，所以 $a_2\in(a_1)$，即 $a_1\mid a_2$。定理证明完毕。
\end{proof}

回忆：对任意环 $R$（不必交换，也不必有 $1$），左 $R$-模 $C$ 称为循环 $R$-模，如果存在 $x\in C$ 使 $C=Rx$。此时可定义 $R$-模同态
\[
\pi:R\longrightarrow C
\]
为 $\pi(r)=rx$；由于 $C=Rx$，这个同态是满射。第一同构定理给出左 $R$-模同构
\[
R/\ker\pi\cong C.
\]
若 $R$ 是 P.I.D.，则 $\ker\pi$ 是主理想 $(a)$，所以循环 $R$-模 $C$ 具有 $R/(a)$ 的形式，其中 $(a)=\operatorname{Ann}(C)$。

循环模是最简单的模（因为只需要一个生成元）。基本定理的存在性部分断言：P.I.D. 上任意有限生成模同构于有限多个循环模的直和。

\noindent\textbf{定理 5（基本定理，存在性：不变因子形式）.} 设 $R$ 是 P.I.D.，$M$ 是有限生成 $R$-模。
\begin{enumerate}[label=(\arabic*)]
\item $M$ 同构于有限多个循环模的直和。更准确地说，
\[
M\cong R^r\oplus R/(a_1)\oplus R/(a_2)\oplus\cdots\oplus R/(a_m)
\]
其中 $r\geq0$ 是整数，$a_1,a_2,\ldots,a_m$ 是 $R$ 中非零且非单位的元素，并满足整除关系
\[
a_1\mid a_2\mid\cdots\mid a_m.
\]
\item $M$ 无挠当且仅当 $M$ 自由。
\item 在 (1) 的分解中，
\[
\operatorname{Tor}(M)\cong R/(a_1)\oplus R/(a_2)\oplus\cdots\oplus R/(a_m).
\]
特别地，$M$ 是挠模当且仅当 $r=0$；在这种情形下，$M$ 的零化子是理想 $(a_m)$。
\end{enumerate}

\begin{proof}
由于 $M$ 按假设可由有限集生成，取 $M$ 的一组生成元 $x_1,x_2,\ldots,x_n$，并要求其基数最小。令 $R^n$ 为秩 $n$ 的自由 $R$-模，基为 $b_1,b_2,\ldots,b_n$。定义同态
\[
\pi:R^n\longrightarrow M
\]
由 $\pi(b_i)=x_i$（对所有 $i$）给出。由于 $x_1,\ldots,x_n$ 生成 $M$，$\pi$ 自动满射。由模的第一同构定理，
\[
R^n/\ker\pi\cong M.
\]
现在对 $R^n$ 及其子模 $\ker\pi$ 应用定理 4，可以选取 $R^n$ 的另一组基 $y_1,y_2,\ldots,y_n$，使得 $a_1y_1,a_2y_2,\ldots,a_my_m$ 是 $\ker\pi$ 的一组基，其中 $a_1,a_2,\ldots,a_m\in R$ 且
\[
a_1\mid a_2\mid\cdots\mid a_m.
\]
这给出
\[
M\cong R^n/\ker\pi=(Ry_1\oplus Ry_2\oplus\cdots\oplus Ry_n)/(Ra_1y_1\oplus Ra_2y_2\oplus\cdots\oplus Ra_my_m).
\]

为了识别右侧商模，我们使用自然满的 $R$-模同态
\[
Ry_1\oplus\cdots\oplus Ry_n\longrightarrow R/(a_1)\oplus R/(a_2)\oplus\cdots\oplus R/(a_m)\oplus R^{n-m}
\]
它把
\[
\alpha_1y_1+\cdots+\alpha_ny_n
\]
送到
\[
(\alpha_1\bmod(a_1),\ldots,\alpha_m\bmod(a_m),\alpha_{m+1},\ldots,\alpha_n).
\]
这个同态的核显然正是那些 $\alpha_i$ 被 $a_i$ 整除（$i=1,2,\ldots,m$）的元素集合，也就是
\[
Ra_1y_1\oplus Ra_2y_2\oplus\cdots\oplus Ra_my_m
\]
（参见练习 7）。于是得到
\[
M\cong R/(a_1)\oplus R/(a_2)\oplus\cdots\oplus R/(a_m)\oplus R^{n-m}.
\]
若某个 $a_i$ 是 $R$ 中单位，则 $R/(a_i)=0$，因此可从直和中删去那些初始的单位因子对应的项。这样得到 (1) 中的分解（其中 $r=n-m$）。

由于对 $R$ 中任意非零元素 $a$，$R/(a)$ 都是挠 $R$-模，(1) 立即蕴含 $M$ 是无挠模当且仅当 $M\cong R^r$，这就是 (2)。而 (3) 也直接由定义得到，因为 $R/(a)$ 的零化子显然是理想 $(a)$，并且整除关系使各循环挠和的总零化子为最大不变因子的理想 $(a_m)$。
\end{proof}

我们很快将证明定理 5 中分解的唯一性：如果还有
\[
M\cong R^{r'}\oplus R/(b_1)\oplus R/(b_2)\oplus\cdots\oplus R/(b_{m'})
\]
其中 $r'\geq0$，$b_1,b_2,\ldots,b_{m'}$ 是非零非单位元素并满足
\[
b_1\mid b_2\mid\cdots\mid b_{m'},
\]
则 $r=r'$，$m=m'$，并且 $(a_i)=(b_i)$（也就是 $a_i$ 与 $b_i$ 只差一个单位因子）对所有 $i$ 成立。正是整除条件 $a_1\mid a_2\mid\cdots\mid a_m$ 给出这种唯一性。

\noindent\textbf{定义.} 定理 5 中的整数 $r$ 称为 $M$ 的\emph{自由秩}或 \emph{Betti 数}；元素 $a_1,a_2,\ldots,a_m\in R$（只差乘以 $R$ 中单位的歧义）称为 $M$ 的\emph{不变因子}。

注意，在证明 $M$ 的不变因子的唯一性之前，严格地说我们应称它们为 $M$ 的一组不变因子（自由秩同理），意思是能给出上述定理 (1) 中分解的任意一组元素。

利用中国剩余定理，可以把定理 5 中的循环模进一步分解，使 $M$ 成为若干循环模的直和，而这些循环模的零化子尽可能简单（即为 $(0)$ 或由 $R$ 中素元幂生成）。这给出另一种分解；我们也会看到它是唯一的。现在描述它。

设 $a$ 是主理想整环 $R$ 中的非零元素。由于 $R$ 也是唯一因子分解整环，可写
\[
a=up_1^{\alpha_1}p_2^{\alpha_2}\cdots p_s^{\alpha_s},
\]
其中 $p_i$ 是 $R$ 中两两不同的素元，$u$ 是单位。这个分解在差一个单位意义下唯一，所以理想 $(p_i^{\alpha_i})$（$i=1,\ldots,s$）唯一确定。若 $i\neq j$，则
\[
(p_i^{\alpha_i})+(p_j^{\alpha_j})=R,
\]
因为这两个理想的和由一个最大公因子生成，而对不同素元 $p_i,p_j$ 来说最大公因子为 $1$。换言之，理想 $(p_i^{\alpha_i})$ 两两互素。所有这些理想的交就是 $(a)$，因为 $a$ 是 $p_1^{\alpha_1},p_2^{\alpha_2},\ldots,p_s^{\alpha_s}$ 的最小公倍元。因此中国剩余定理（定理 7.17）说明
\[
R/(a)\cong R/(p_1^{\alpha_1})\oplus R/(p_2^{\alpha_2})\oplus\cdots\oplus R/(p_s^{\alpha_s})
\]
既是环同构，也是 $R$-模同构。

把这应用到定理 5 中的模，可将 $M$ 的每个不变因子 $a_i$ 对应的直和项 $R/(a_i)$ 写成若干循环模的直和，这些循环模的零化子为 $a_i$ 的素幂因子。于是得到：

\noindent\textbf{定理 6（基本定理，存在性：初等因子形式）.} 设 $R$ 是 P.I.D.，$M$ 是有限生成 $R$-模。则 $M$ 是有限多个循环模的直和，其零化子或者为 $(0)$，或者由 $R$ 中素元的幂生成；也就是说，
\[
M\cong R^r\oplus R/(p_1^{\alpha_1})\oplus R/(p_2^{\alpha_2})\oplus\cdots\oplus R/(p_t^{\alpha_t}),
\]
其中 $r\geq0$ 为整数，$p_1^{\alpha_1},\ldots,p_t^{\alpha_t}$ 是 $R$ 中（不必互异的）素元的正幂。

我们是用 $M$ 的不变因子的素幂因子证明定理 6 的。事实上，我们将看到，把 $M$ 分解为定理 6 中零化子为 $(0)$ 或素幂的循环模直和也是唯一的；也就是说，整数 $r$ 与理想 $(p_1^{\alpha_1}),\ldots,(p_t^{\alpha_t})$ 由 $M$ 唯一确定。这些素幂有名称：

\noindent\textbf{定义.} 设 $R$ 是 P.I.D.，$M$ 是如定理 6 中的有限生成 $R$-模。素幂 $p_1^{\alpha_1},\ldots,p_t^{\alpha_t}$（只差乘以 $R$ 中单位的歧义）称为 $M$ 的\emph{初等因子}。

设 $M$ 是主理想整环 $R$ 上的有限生成挠模。如果对定理 6 的分解中出现的不同素元 $p_1,p_2,\ldots,p_n$，把所有对应同一素元 $p_i$ 的循环因子聚在一起，我们特别得到 $M$ 可写成直和
\[
M=N_1\oplus N_2\oplus\cdots\oplus N_n,
\]
其中 $N_i$ 由 $M$ 中被素元 $p_i$ 的某个幂零化的所有元素组成。这个结论对不一定有限生成的 $R$-模也成立：

\noindent\textbf{定理 7（Primary 分解定理）.} 设 $R$ 是 P.I.D.，$M$ 是非零挠 $R$-模（不必有限生成），且具有非零零化子 $a$。设 $a$ 分解为 $R$ 中不同素元的幂的乘积：
\[
a=up_1^{\alpha_1}p_2^{\alpha_2}\cdots p_n^{\alpha_n},
\]
并令
\[
N_i=\{x\in M\mid p_i^{\alpha_i}x=0\},\qquad 1\leq i\leq n.
\]
则 $N_i$ 是 $M$ 的子模，零化子为 $(p_i^{\alpha_i})$，并且它是所有被 $p_i$ 的某个幂零化的元素构成的子模。我们有
\[
M=N_1\oplus N_2\oplus\cdots\oplus N_n.
\]
如果 $M$ 有限生成，则每个 $N_i$ 是有限多个循环模的直和，其零化子均为 $p_i^{\alpha_i}$ 的因子。

\begin{proof}
当 $M$ 在 $R$ 上有限生成时，这些结论已经由定理 6 及其后的分组讨论得到。在一般情形下，显然每个 $N_i$ 是 $M$ 的子模，且其零化子整除 $(p_i^{\alpha_i})$。由于 $R$ 是 P.I.D.，当 $i\neq j$ 时理想 $(p_i^{\alpha_i})$ 与 $(p_j^{\alpha_j})$ 互素，因此可以把中国剩余定理证明中的论证轻微修改后应用于模，得到 $M$ 的这个直和分解。利用这个直和分解，容易看出 $N_i$ 的零化子恰为 $(p_i^{\alpha_i})$。
\end{proof}

\noindent\textbf{定义.} 前一定理中的子模 $N_i$ 称为 $M$ 的 $p_i$-\emph{primary 分量}。

用这个术语来说，有限生成模 $M$ 的初等因子正是 $\operatorname{Tor}(M)$ 的 primary 分量的不变因子。

现在证明关于基本定理中分解的唯一性陈述。注意，如果 $M$ 是交换环 $R$ 上的任意模，$a\in R$，则
\[
aM=\{am\mid m\in M\}
\]
是 $M$ 的子模。还回忆，在主理想整环 $R$ 中，非零素理想都是极大理想，所以 $R$ 对非零素理想的商是域。

\noindent\textbf{引理 8.} 设 $R$ 是 P.I.D.，$p$ 是 $R$ 中素元，$F=R/(p)$。
\begin{enumerate}[label=(\arabic*)]
\item 若 $M=R^r$，则 $M/pM\cong F^r$。
\item 若 $M=R/(a)$，其中 $a\in R$ 非零，则
\[
M/pM\cong
\begin{cases}
F,& p\mid a,\\
0,& p\nmid a.
\end{cases}
\]
\item 若 $M=R/(a_1)\oplus R/(a_2)\oplus\cdots\oplus R/(a_k)$，且每个 $a_i$ 都被 $p$ 整除，则 $M/pM\cong F^k$。
\end{enumerate}

\begin{proof}
(1) 从 $R^r$ 到 $(R/(p))^r$ 有自然映射，把 $(a_1,\ldots,a_r)$ 送到 $(a_1\bmod(p),\ldots,a_r\bmod(p))$。这显然是满的 $R$-模同态，它的核由所有坐标都被 $p$ 整除的 $r$ 元组组成，也就是 $pR^r$。故 $R^r/pR^r\cong(R/(p))^r=F^r$。

(2) 这由同构定理得到。先注意 $p(R/(a))$ 是理想 $(p)$ 在商 $R/(a)$ 中的像，因此等于
\[
((p)+(a))/(a).
\]
理想 $(p)+(a)$ 由 $p$ 与 $a$ 的一个最大公因子生成；所以当 $p\mid a$ 时它是 $(p)$，当 $p\nmid a$ 时它是 $R=(1)$。因此当 $p\mid a$ 时 $pM=(p)/(a)$，当 $p\nmid a$ 时 $pM=R/(a)=M$。若 $p\mid a$，则
\[
M/pM=(R/(a))/((p)/(a))\cong R/(p)=F;
\]
若 $p\nmid a$，则 $M/pM=M/M=0$。

(3) 由 (2) 对直和逐项应用得到，正如定理 5(1) 证明中识别直和商的方式一样。
\end{proof}

\noindent\textbf{定理 9（基本定理，唯一性）.} 设 $R$ 是 P.I.D.。
\begin{enumerate}[label=(\arabic*)]
\item 两个有限生成 $R$-模 $M_1$ 与 $M_2$ 同构，当且仅当它们有相同的自由秩和相同的不变因子列表。
\item 两个有限生成 $R$-模 $M_1$ 与 $M_2$ 同构，当且仅当它们有相同的自由秩和相同的初等因子列表。
\end{enumerate}

\begin{proof}
如果 $M_1$ 与 $M_2$ 有相同的自由秩和相同的不变因子列表，或者有相同的自由秩和相同的初等因子列表，则它们显然同构。

反过来，设 $M_1\cong M_2$。任一 $M_1$ 与 $M_2$ 之间的同构都把 $M_1$ 的挠部分映到 $M_2$ 的挠部分，所以
\[
\operatorname{Tor}(M_1)\cong\operatorname{Tor}(M_2).
\]
于是
\[
R^{r_1}\cong M_1/\operatorname{Tor}(M_1)
\cong M_2/\operatorname{Tor}(M_2)\cong R^{r_2},
\]
其中 $r_1$、$r_2$ 分别是 $M_1$、$M_2$ 的自由秩。令 $p$ 为 $R$ 中任意非零素元。由 $R^{r_1}\cong R^{r_2}$ 可得
\[
R^{r_1}/pR^{r_1}\cong R^{r_2}/pR^{r_2}.
\]
由引理 8(1)，这说明 $F^{r_1}\cong F^{r_2}$，其中 $F=R/pR$ 是域。因此两个 $F$-向量空间同构，维数相等，故 $r_1=r_2$。所以 $M_1$ 与 $M_2$ 有相同的自由秩。

于是只需证明 $M_1$ 与 $M_2$ 有相同的不变因子列表和初等因子列表。为此只需要处理同构的挠模 $\operatorname{Tor}(M_1)$ 与 $\operatorname{Tor}(M_2)$；也就是说，不妨假设 $M_1,M_2$ 都是挠 $R$-模。

先证明它们有相同的初等因子。只需说明：对任一固定素元 $p$，作为 $p$ 的幂的初等因子在 $M_1$ 与 $M_2$ 中相同。如果 $M_1\cong M_2$，则 $M_1$ 的 $p$-primary 子模（即初等因子为 $p$ 的幂的循环因子的直和）同构于 $M_2$ 的 $p$-primary 子模，因为这些子模正是所有被 $p$ 的某个幂零化的元素组成的子模。因此归结为证明：若两个零化子为 $p$ 的幂的模 $M_1,M_2$ 同构，则它们有相同的初等因子。

对 $M_1$ 的零化子中 $p$ 的幂指数作归纳（这与 $M_2$ 中的指数相同，因为两者同构）。若该幂指数为 $0$，则 $M_1=M_2=0$，结论成立。否则 $M_1$ 与 $M_2$ 都有非平凡初等因子。设 $M_1$ 的初等因子为
\[
\underbrace{p,p,\ldots,p}_{m\text{ 次}},\quad p^{\alpha_1},p^{\alpha_2},\ldots,p^{\alpha_s},
\]
其中 $2\leq\alpha_1\leq\alpha_2\leq\cdots\leq\alpha_s$；也就是说，$M_1$ 是若干循环模的直和，生成元可记为 $x_1,x_2,\ldots,x_m,x_{m+1},\ldots,x_{m+s}$，其零化子分别为 $(p),\ldots,(p),(p^{\alpha_1}),\ldots,(p^{\alpha_s})$。那么子模 $pM_1$ 的初等因子为
\[
p^{\alpha_1-1},p^{\alpha_2-1},\ldots,p^{\alpha_s-1},
\]
因为 $pM_1$ 是以 $px_1,\ldots,px_m,px_{m+1},\ldots,px_{m+s}$ 生成的循环模直和，而这些生成元的零化子为 $(1),\ldots,(1),(p^{\alpha_1-1}),\ldots,(p^{\alpha_s-1})$。

类似地，设 $M_2$ 的初等因子为
\[
\underbrace{p,p,\ldots,p}_{n\text{ 次}},\quad p^{\beta_1},p^{\beta_2},\ldots,p^{\beta_t},
\]
其中 $2\leq\beta_1\leq\beta_2\leq\cdots\leq\beta_t$，则 $pM_2$ 的初等因子为
\[
p^{\beta_1-1},p^{\beta_2-1},\ldots,p^{\beta_t-1}.
\]
由于 $M_1\cong M_2$，也有 $pM_1\cong pM_2$；并且 $pM_1$ 的零化子中 $p$ 的幂指数比 $M_1$ 的少 $1$。由归纳假设，$pM_1$ 与 $pM_2$ 的初等因子相同，于是 $s=t$ 且 $\alpha_i-1=\beta_i-1$，即 $\alpha_i=\beta_i$（$1\leq i\leq s$）。最后，因为
\[
M_1/pM_1\cong M_2/pM_2,
\]
由引理 8(3) 得到 $F^{m+s}\cong F^{n+t}$，从而 $m+s=n+t$；已知 $s=t$，故 $m=n$。这证明了 $M_1$ 与 $M_2$ 的初等因子集合相同。

现在证明 $M_1$ 与 $M_2$ 必有相同的不变因子。设
\[
a_1\mid a_2\mid\cdots\mid a_m
\]
是 $M_1$ 的不变因子。把这些元素分解为素幂因子，得到 $M_1$ 的一组初等因子。注意整除关系意味着 $a_m$ 是这些初等因子中每个不同素元所对应的最大素幂的乘积；去掉 $a_m$ 的这些因子后，$a_{m-1}$ 是剩余初等因子中每个不同素元所对应最大素幂的乘积；依此类推。若
\[
b_1\mid b_2\mid\cdots\mid b_n
\]
是 $M_2$ 的不变因子，同样地通过分解得到 $M_2$ 的一组初等因子。上面已经证明 $M_1$ 与 $M_2$ 的初等因子相同，因此由刚才的重组规则可知它们的不变因子也相同。
\end{proof}

\noindent\textbf{推论 10.} 设 $R$ 是 P.I.D.，$M$ 是有限生成 $R$-模。
\begin{enumerate}[label=(\arabic*)]
\item $M$ 的初等因子是 $M$ 的不变因子的素幂因子。
\item $M$ 的最大不变因子是 $M$ 的初等因子中互异素元对应的最大素幂的乘积；次大的不变因子是剩余初等因子中互异素元对应最大素幂的乘积；依此类推。
\end{enumerate}

\begin{proof}
(1) 中的过程给出一组初等因子，而初等因子由定理 9 唯一，所以 (1) 中的过程给出的就是初等因子的集合。(2) 同理。
\end{proof}

\noindent\textbf{推论 11（有限生成阿贝尔群基本定理）.} 见定理 5.3 与定理 5.5。

\begin{proof}
在定理 5、6 与 9 中取 $R=\mathbb Z$ 即可（不过注意，第 5 章中为了计算方便，不变因子的列出顺序与这里相反）。
\end{proof}

推论 10 中在初等因子与不变因子之间转换的过程，在第 5 章有限生成阿贝尔群的情形中有较详细的描述。还应注意，如果一个有限生成模 $M$ 写成形如 $R/(a)$ 的循环模的直和，则出现的理想 $(a)$ 一般并不唯一，除非附加某些条件（例如不变因子的整除条件，或者在初等因子情形中要求 $a$ 是素元的幂）。要判断两个模是否同构，必须先把它们写成这种标准（或规范）形式。
\subsection*{练习}

\begin{enumerate}[label=\arabic*.]
\item 设 $M$ 是整环 $R$ 上的模。
\begin{enumerate}[label=(\alph*)]
\item 设 $x$ 是 $M$ 中非零挠元素。证明 $x$ 与 $0$ ``线性相关''。推出 $\operatorname{Tor}(M)$ 的秩为 $0$，因而任意挠 $R$-模的秩为 $0$。
\item 证明 $M$ 的秩等于无挠商模 $M/\operatorname{Tor}M$ 的秩。
\end{enumerate}

\item 设 $M$ 是整环 $R$ 上的模。
\begin{enumerate}[label=(\alph*)]
\item 设 $M$ 的秩为 $n$，且 $x_1,x_2,\ldots,x_n$ 是 $M$ 中任意极大线性无关集。令 $N=Rx_1+\cdots+Rx_n$。证明 $N\cong R^n$，且 $M/N$ 是挠 $R$-模。等价地，$x_1,\ldots,x_n$ 线性无关，并且对任意 $y\in M$，存在非零 $r\in R$，使得 $ry$ 可写为 $r_1x_1+\cdots+r_nx_n$。
\item 反过来，若 $M$ 含有自由秩为 $n$ 的子模 $N$（即 $N\cong R^n$），并且 $M/N$ 是挠 $R$-模，证明 $M$ 的秩为 $n$。提示：令 $y_1,\ldots,y_{n+1}$ 为 $M$ 中任意 $n+1$ 个元素。利用 $M/N$ 为挠模，取非零 $r_i\in R$ 使 $r_i y_i$ 都可写成 $N$ 的一组基的线性组合，再仿照命题 3 的证明说明 $r_i y_i$，从而 $y_i$，线性相关。
\end{enumerate}

\item 设 $R$ 为整环，$A,B$ 分别为秩 $m,n$ 的 $R$-模。证明 $A\oplus B$ 的秩为 $m+n$。提示：使用上一题。

\item 设 $R$ 为整环，$M$ 为 $R$-模，$N$ 为 $M$ 的子模。设 $M$ 的秩为 $n$，$N$ 的秩为 $r$，商模 $M/N$ 的秩为 $s$。证明 $n=r+s$。提示：取 $x_1,\ldots,x_s\in M$，其像在 $M/N$ 中构成极大线性无关集；再取 $x_{s+1},\ldots,x_{s+r}$ 为 $N$ 中极大线性无关集。证明这些元素在 $M$ 中线性无关，并且任意 $y\in M$ 的某个非零倍数可写成它们的线性组合，然后使用练习 2。

\item 令 $R=\mathbb Z[x]$，令 $M=(2,x)$ 为由 $2$ 与 $x$ 生成的理想，把它看作 $R$ 的子模。证明 $\{2,x\}$ 不是 $M$ 的基。提示：找出这两个元素之间的一个非平凡 $R$-线性相关关系。证明 $M$ 的秩为 $1$，但 $M$ 不是秩 $1$ 的自由模（参见练习 2）。

\item 证明：若 $R$ 是整环，$M$ 是 $R$ 的任意非主理想，则 $M$ 是无挠、秩 $1$ 的 $R$-模，但不是自由 $R$-模。

\item 设 $R$ 为任意环，$A_1,A_2,\ldots,A_m$ 为 $R$-模，$B_i$ 为 $A_i$ 的子模，$1\leq i\leq m$。证明
\[
(A_1\oplus\cdots\oplus A_m)/(B_1\oplus\cdots\oplus B_m)
\cong (A_1/B_1)\oplus\cdots\oplus(A_m/B_m).
\]

\item 设 $R$ 为 P.I.D.，$B$ 为挠 $R$-模，$p$ 为 $R$ 中素元。证明：若存在非零 $b\in B$ 满足 $pb=0$，则 $\operatorname{Ann}(B)\subseteq (p)$。

\item 给出一个整环 $R$ 与非零挠 $R$-模 $M$ 的例子，使得 $\operatorname{Ann}(M)=0$。证明：若 $N$ 是有限生成挠 $R$-模，则 $\operatorname{Ann}(N)\ne0$。

\item 设 $p$ 为 P.I.D. $R$ 中的素元，$N$ 为 $R$-模。证明 $N$ 的 $p$-primary 分量是 $N$ 的子模，并证明 $N$ 是其各 $p$-primary 分量的直和（分量个数不必有限）。

\item 设 $R$ 为 P.I.D.，$a\in R$ 非零，$M=R/(a)$。对 $R$ 的任意素元 $p$，证明
\[
p^{k-1}M/p^kM\cong
\begin{cases}
R/(p),& 1\leq k\leq n,\\
0,& k>n,
\end{cases}
\]
其中 $n$ 是 $a$ 中 $p$ 的幂指数。

\item 设 $R$ 为 P.I.D.，$p$ 为 $R$ 中素元。
\begin{enumerate}[label=(\alph*)]
\item 设 $M$ 是有限生成挠 $R$-模。用上一题证明
\[
p^{k-1}M/p^kM\cong F^{n_k},
\]
其中 $F=R/(p)$，$n_k$ 是 $M$ 的初等因子中形如 $p^a$ 且 $a\geq k$ 的个数。
\item 设 $M_1,M_2$ 是同构的有限生成挠 $R$-模。用 (a) 证明：对每个 $k\geq0$，$M_1$ 与 $M_2$ 有相同数目的初等因子 $p^a$，其中 $a\geq k$。证明这推出 $M_1$ 与 $M_2$ 有相同的初等因子集合。
\end{enumerate}

\item 若 $M$ 是 P.I.D. $R$ 上的有限生成模，描述 $M/\operatorname{Tor}(M)$ 的结构。

\item 设 $R$ 为 P.I.D.，$M$ 为挠 $R$-模。证明 $M$ 不可约（参见第 10.3 节练习 9--11）当且仅当对每个非零 $m\in M$，都有 $M=Rm$，且 $m$ 的零化子是某个非零素理想 $(p)$。

\item 证明：若 $R$ 是 Noether 环，则 $R^n$ 是 Noether $R$-模。提示：固定 $R^n$ 的一组基。若 $M$ 是 $R^n$ 的子模，说明 $M$ 中元素的第一坐标构成 $R$ 的子模，因而有限生成。取 $m_1,
\ldots,m_k\in M$，使它们的第一坐标生成该子模。证明 $M$ 的任意元素都可写成这些 $m_i$ 的 $R$-线性组合加上一个第一坐标为 $0$ 的 $M$ 中元素。再证明 $M\cap R^{n-1}$ 是 $R^{n-1}$ 的子模，其中 $R^{n-1}$ 表示 $R^n$ 中第一坐标为 $0$ 的元素集合，最后对 $n$ 归纳。

\medskip
以下一组练习给出定理 5 在 $R$ 为 Euclidean 整环时的矩阵证明，使用行、列初等变换。这尤其适用于 $R=\mathbb Z$ 与 $R=F[x]$，并且在计算上很有用。设 $R$ 为 Euclidean 整环，$M$ 为 $R$-模。

\item 证明 $M$ 有限生成当且仅当存在某个整数 $n$ 和满的 $R$-同态 $\varphi:R^n\to M$。（这对任意环 $R$ 都成立。）

设 $\varphi:R^n\to M$ 为满 $R$-模同态。由练习 15，$\ker\varphi$ 有限生成。若 $x_1,
\ldots,x_n$ 是 $R^n$ 的基，$y_1,
\ldots,y_m$ 是 $\ker\varphi$ 的生成元，则
\[
y_i=a_{i1}x_1+a_{i2}x_2+\cdots+a_{in}x_n,
\qquad i=1,2,\ldots,m,
\]
其中 $a_{ij}\in R$。因此同态 $\varphi$（也就是 $M$ 的模结构）由 $R^n$ 的生成元选择和矩阵 $A=(a_{ij})$ 决定。这样的矩阵称为\emph{关系矩阵}。

\item
\begin{enumerate}[label=(\alph*)]
\item 证明在 $R^n$ 的基中交换 $x_i$ 与 $x_j$，会交换相应关系矩阵的第 $i$ 列和第 $j$ 列。
\item 证明：对任意 $a\in R$，把基元素 $x_j$ 替换为 $x_j-ax_i$ 后仍得到 $R^n$ 的一组基；相应关系矩阵等于原关系矩阵，只是把 $a$ 倍的第 $j$ 列加到第 $i$ 列上。提示：
\[
\cdots+a_ix_i+\cdots+a_jx_j+\cdots
=\cdots+(a_i+aa_j)x_i+\cdots+a_j(x_j-ax_i)+\cdots.
\]
\end{enumerate}

\item
\begin{enumerate}[label=(\alph*)]
\item 证明交换生成元 $y_i,y_j$ 会交换关系矩阵的第 $i$ 行和第 $j$ 行。
\item 证明：对任意 $a\in R$，把 $y_j$ 替换为 $y_j-ay_i$ 仍给出 $\ker\varphi$ 的一组生成元；相应关系矩阵等于原矩阵，只是把 $-a$ 倍的第 $i$ 行加到第 $j$ 行上。
\end{enumerate}

\item 由前两题可知，可通过重新选择 $R^n$ 与 $\ker\varphi$ 的生成元，对给定关系矩阵作初等行、列变换。若所有关系矩阵都是零矩阵，则 $\ker\varphi=0$，$M\cong R^n$。否则，令 $a_1$ 为某个初始关系矩阵中所有非零元素的最大公因子。
\begin{enumerate}[label=(\alph*)]
\item 证明可用初等行、列变换使 $a_1$ 出现在关系矩阵左上角，并且所得矩阵中 $a_1$ 整除所有元素。
\item 证明存在如下形式的关系矩阵
\[
\begin{pmatrix}
 a_1&0&\cdots&0\\
 0&a_{22}&\cdots&a_{2n}\\
 \vdots&\vdots&\ddots&\vdots\\
 0&a_{m2}&\cdots&a_{mn}
\end{pmatrix},
\]
其中 $a_1$ 整除所有条目。
\item 令 $a_2$ 为 (b) 中除 $a_1$ 之外所有条目的最大公因子。证明存在如下形式的关系矩阵
\[
\begin{pmatrix}
 a_1&0&0&\cdots&0\\
 0&a_2&0&\cdots&0\\
 0&0&a_{33}&\cdots&a_{3n}\\
 \vdots&\vdots&\vdots&\ddots&\vdots
\end{pmatrix},
\]
其中 $a_1\mid a_2$，且 $a_2$ 整除其余所有条目。
\item 证明存在形如
\[
\begin{pmatrix}
D&0\\
0&0
\end{pmatrix}
\]
的关系矩阵，其中 $D$ 为对角矩阵，非零对角元为 $a_1,a_2,\ldots,a_k$，$k\leq n$，并满足
\[
a_1\mid a_2\mid\cdots\mid a_k.
\]
推出
\[
M\cong R/(a_1)\oplus R/(a_2)\oplus\cdots\oplus R/(a_k)\oplus R^{n-k}.
\]
\end{enumerate}

若 $n$ 不是生成 $M$ 所需的最小生成元个数，则上述某些初始 $a_i$ 可能是单位；对应直和项为零。去掉这些无关因子即可得到 $M$ 的不变因子。并且，在上面的约化过程中跟踪列变换，就可得到不变因子分解中各循环子模的 $R$-生成元。

\item 设 $R$ 为整环，$F$ 为其分式域，$M$ 为任意 $R$-模。证明 $M$ 的秩等于 $F\otimes_R M$ 作为 $F$-向量空间的维数。

\item 证明 P.I.D. 上的有限生成模是投射模当且仅当它是自由模。

\item 设 $R$ 为不是域的 P.I.D.。证明没有有限生成 $R$-模是内射模。提示：用第 10.5 节练习 4 分别考虑挠模与自由模。
\end{enumerate}



\section{有理标准形}

现在把上一节的结果应用到 P.I.D. 为 $F[x]$ 的特殊情形，其中 $F[x]$ 是域 $F$ 上关于 $x$ 的多项式环。

令 $V$ 是 $F$ 上维数为 $n$ 的有限维向量空间，$T$ 是 $V$ 的一个固定线性变换（即 $V$ 到自身的线性变换）。正如第 10 章所见，可以把 $V$ 看作 $F[x]$-模，其中元素 $x$ 在 $V$ 上按线性变换 $T$ 作用（于是 $x$ 中任意多项式按 $T$ 中相同多项式作用）。由于 $V$ 按假设是 $F$ 上有限维的，所以按定义它是有限生成 $F$-模，因而当然是有限生成 $F[x]$-模；所以前一节的分类定理适用。

任何非零自由 $F[x]$-模（同构于若干个 $F[x]$ 的直和）作为 $F$-向量空间都是无限维的，因此如果 $V$ 在 $F$ 上有限维，那么它作为 $F[x]$-模事实上必为挠模（即自由秩为 $0$）。于是由基本定理，$V$ 作为 $F[x]$-模同构于有限多个循环挠 $F[x]$-模的直和。我们将看到，这个 $V$ 的分解允许我们选取 $V$ 的一组基，使得线性变换 $T$ 的矩阵表示具有一种特定的简单形式。当使用 $V$ 的不变因子分解时，得到 $T$ 的矩阵的有理标准形；本节分析这种形式。当使用初等因子分解（并且 $F$ 含有 $T$ 的所有特征值）时，得到下一节讨论的 Jordan 标准形，也就是前面提到的尽可能接近对角矩阵的 $T$ 的矩阵。基本定理的唯一性部分保证有理标准形和 Jordan 标准形是唯一的（因此称为标准形）。

这些标准形的一个重要用途，是分类 $V$ 的不同线性变换。特别地，它们使我们能够判定两个矩阵何时表示同一个线性变换，即判定给定的两个 $n\times n$ 矩阵何时相似。

注意，这又是被作用空间的结构（例如 $V$ 的不变因子分解）用来获得作用在其上的代数对象（这里是线性变换）的重要信息的一个例子。第 18 章将在群作用在向量空间上的情形中讨论这一点（这称为群的表示论）。

在详细描述有理标准形之前，我们先引入一些线性代数概念。

\noindent\textbf{定义.}
\begin{enumerate}[label=(\arabic*)]
\item 如果存在非零向量 $v\in V$ 使 $T(v)=\lambda v$，则 $F$ 中元素 $\lambda$ 称为线性变换 $T$ 的\emph{特征值}。在这种情形下，$v$ 称为 $T$ 的一个对应于特征值 $\lambda$ 的\emph{特征向量}。
\item 如果 $A$ 是一个系数在 $F$ 中的 $n\times n$ 矩阵，$\lambda\in F$，并且存在非零 $n\times1$ 列向量 $v$ 使 $Av=\lambda v$，则称 $\lambda$ 为 $A$ 的特征值，$v$ 为对应特征向量。
\item 如果 $\lambda$ 是线性变换 $T$ 的特征值，则集合
\[
\{v\in V\mid T(v)=\lambda v\}
\]
称为 $T$ 对应于特征值 $\lambda$ 的\emph{特征空间}。类似地，如果 $\lambda$ 是 $n\times n$ 矩阵 $A$ 的特征值，则满足 $Av=\lambda v$ 的 $n\times1$ 矩阵 $v$ 的集合称为 $A$ 对应于特征值 $\lambda$ 的特征空间。
\end{enumerate}

注意，如果固定 $V$ 的一组基 $\mathcal B$，则 $V$ 的任意线性变换 $T$ 都有一个对应的 $n\times n$ 矩阵 $A$。反过来，若 $A$ 是任意 $n\times n$ 矩阵，则由 $T(v)=Av$ 定义的映射（右边的 $v$ 指 $v$ 在固定基 $\mathcal B$ 下的坐标列向量）是 $V$ 的线性变换。于是 $v$ 是 $T$ 的特征值 $\lambda$ 的特征向量，当且仅当 $v$ 在 $\mathcal B$ 下的坐标列向量是 $A$ 的特征值 $\lambda$ 的特征向量。换言之，线性变换 $T$ 的特征值与 $T$ 在任意固定基下矩阵 $A$ 的特征值相同。

\noindent\textbf{定义.} $V$ 到自身的线性变换的行列式定义为表示该线性变换的任意矩阵的行列式（注意这与所选基无关）。

\noindent\textbf{命题 12.} 下列条件等价：
\begin{enumerate}[label=(\arabic*)]
\item $\lambda$ 是 $T$ 的特征值；
\item $\lambda I-T$ 是 $V$ 的奇异线性变换；
\item $\det(\lambda I-T)=0$。
\end{enumerate}

\begin{proof}
$\lambda$ 是 $T$ 的特征值且有对应特征向量 $v$，当且仅当 $v$ 是 $\lambda I-T$ 的核中的非零向量。因此 (1) 与 (2) 等价。由关于行列式的结果，线性变换奇异当且仅当其行列式为零，故 (2) 与 (3) 等价。
\end{proof}

\noindent\textbf{定义.} 令 $x$ 为 $F$ 上的不定元。多项式
\[
\det(xI-T)
\]
称为 $T$ 的\emph{特征多项式}，记为 $c_T(x)$。如果 $A$ 是系数在 $F$ 中的 $n\times n$ 矩阵，则 $\det(xI-A)$ 称为 $A$ 的特征多项式，记为 $c_A(x)$。

由展开行列式容易看出，$T$ 或 $A$ 的特征多项式都是次数 $n=\dim V$ 的首一多项式。命题 12 说明 $T$（或 $A$）的特征值集合正是 $T$（或 $A$）特征多项式的根的集合。特别地，$T$ 至多有 $n$ 个不同的特征值。

我们已经看到，$V$ 通过线性变换 $T$ 看作 $F[x]$-模时是挠 $F[x]$-模。令 $m(x)\in F[x]$ 是生成 $F[x]$ 中 $\operatorname{Ann}(V)$ 的唯一首一多项式。等价地，$m(x)$ 是使 $V$ 被零化的最低次数首一多项式（即 $m(T)$ 是零线性变换），并且若 $f(x)\in F[x]$ 是任意零化 $V$ 的多项式，则 $m(x)$ 整除 $f(x)$。由于选择 $V$ 的一组基可给出所有 $n\times n$ 矩阵环与 $V$ 到自身线性变换集合的同构，对 $F$ 上任意 $n\times n$ 矩阵 $A$，同样存在唯一首一最低次数多项式 $m(x)$ 使 $m(A)$ 为零矩阵。

\noindent\textbf{定义.} 生成 $F[x]$ 中理想 $\operatorname{Ann}(V)$ 的唯一首一多项式称为 $T$ 的\emph{最小多项式}，记为 $m_T(x)$。在矩阵 $A$ 上取值为零矩阵的最低次数首一多项式称为 $A$ 的\emph{最小多项式}，记为 $m_A(x)$。

容易看出（参见练习 5），这些最小多项式的次数至多为 $n^2$，其中 $n=\dim V$。我们很快将证明 $T$ 的最小多项式整除 $T$ 的特征多项式（这是 Cayley--Hamilton 定理），矩阵情形同理；所以实际上这些多项式的次数至多为 $n$。

现在描述线性变换 $T$（以及矩阵 $A$）的有理标准形。由定理 5，有 $F[x]$-模同构
\[
V\cong F[x]/(a_1(x))\oplus F[x]/(a_2(x))\oplus\cdots\oplus F[x]/(a_m(x)),\tag{12.1}
\]
其中 $a_1(x),a_2(x),\ldots,a_m(x)$ 是 $F[x]$ 中次数至少为 $1$ 的多项式，并满足整除条件
\[
a_1(x)\mid a_2(x)\mid\cdots\mid a_m(x).
\]
这些不变因子只在差一个 $F[x]$ 中单位的意义下确定；由于 $F[x]$ 的单位恰是 $F$ 的非零元素（非零常数多项式），我们可规定这些多项式为首一，从而使它们唯一。

由于 $V$ 的零化子是理想 $(a_m(x))$（定理 5(3)），立即得到：

\noindent\textbf{命题 13.} 最小多项式 $m_T(x)$ 是 $V$ 的最大不变因子。$V$ 的所有不变因子都整除 $m_T(x)$。

我们下面将看到如何计算不仅 $T$ 的最小多项式，也计算其他不变因子。

现在为 (12.1) 中 $V$ 的每个直和因子选取一组基，使得 $T$ 的矩阵非常简单。回忆：等式 (12.1) 左侧的线性变换 $T$ 就是右侧每个因子上元素 $x$ 的乘法作用。

在第 11 章命题 1 后的例子中已经看到，若
\[
a(x)=x^k+b_{k-1}x^{k-1}+\cdots+b_1x+b_0
\]
是 $F[x]$ 中任意首一多项式，并令 $\bar x=x\bmod(a(x))$，则元素
\[
1,\bar x,\bar x^2,\ldots,\bar x^{k-1}
\]
给出 $F[x]/(a(x))$ 的 $F$-向量空间基。关于这组基，乘以 $x$ 的线性变换按如下方式作用：
\[
1\mapsto\bar x,
\quad \bar x\mapsto\bar x^2,
\quad\ldots,
\quad \bar x^{k-2}\mapsto\bar x^{k-1},
\]
而
\[
\bar x^{k-1}\mapsto \bar x^k=-b_0-b_1\bar x-\cdots-b_{k-1}\bar x^{k-1},
\]
最后一个等式来自 $a(\bar x)=0$。因此关于这组基，乘以 $x$ 的矩阵是
\[
\begin{pmatrix}
0&0&0&\cdots&0&-b_0\\
1&0&0&\cdots&0&-b_1\\
0&1&0&\cdots&0&-b_2\\
\vdots&\vdots&\vdots& &\vdots&\vdots\\
0&0&0&\cdots&1&-b_{k-1}
\end{pmatrix}.
\]

\noindent\textbf{定义.} 上述矩阵称为首一多项式 $a(x)$ 的\emph{伴随矩阵}，记为 $C(a)$。

\noindent\textbf{定义.} 若一个块对角矩阵的对角块为首一多项式 $a_1(x),\ldots,a_m(x)$ 的伴随矩阵 $C(a_1),\ldots,C(a_m)$，并且
\[
a_1(x)\mid a_2(x)\mid\cdots\mid a_m(x),
\]
则称该矩阵处于\emph{有理标准形}。

\noindent\textbf{定理 14（线性变换的有理标准形）.} 设 $V$ 是域 $F$ 上有限维向量空间，$T$ 是 $V$ 上线性变换。则存在 $V$ 的一组基，使得 $T$ 关于这组基的矩阵处于有理标准形。这个有理标准形由 $V$ 作为 $F[x]$-模的不变因子唯一确定。

\begin{proof}
把 $V$ 看作由 $T$ 定义的 $F[x]$-模。由定理 5，$V$ 分解为循环挠模 $F[x]/(a_i(x))$ 的直和，其中 $a_1\mid\cdots\mid a_m$。在每个因子 $F[x]/(a_i(x))$ 上取基 $1,\bar x,\ldots,\bar x^{d_i-1}$。上面的计算说明，在这一基上元素 $x$ 的作用（即 $T$ 的作用）由伴随矩阵 $C(a_i)$ 给出。把这些基依直和顺序合在一起，得到 $V$ 的一组基；在这组基下 $T$ 的矩阵是这些伴随矩阵的块对角直和，因而处于有理标准形。唯一性来自定理 9 中不变因子的唯一性。
\end{proof}

\noindent\textbf{定理 15.} 设 $S,T$ 是 $V$ 的线性变换。下列条件等价：
\begin{enumerate}[label=(\arabic*)]
\item $S$ 与 $T$ 相似，即存在 $V$ 的一个非奇异线性变换 $P$，使得 $S=P^{-1}TP$；
\item $S$ 与 $T$ 使 $V$ 成为同构的 $F[x]$-模；
\item $S$ 与 $T$ 有相同的不变因子。
\end{enumerate}

\begin{proof}
若 $S=P^{-1}TP$，则 $P$ 给出两个 $F[x]$-模之间的同构：在第一个模中 $x$ 按 $S$ 作用，在第二个模中 $x$ 按 $T$ 作用；等式 $TP=PS$ 正是 $P$ 与 $x$ 的作用可交换。反过来，若存在 $F[x]$-模同构 $P$，则 $P$ 作为 $F$-线性映射非奇异，并满足 $P(Sv)=T(Pv)$，故 $S=P^{-1}TP$。条件 (2) 与 (3) 的等价性正是定理 9 对 P.I.D. $F[x]$ 的应用。
\end{proof}

\noindent\textbf{定理 16（矩阵的有理标准形）.} 设 $A$ 是域 $F$ 上的 $n\times n$ 矩阵。则 $A$ 相似于一个有理标准形矩阵，也就是说，存在 $F$ 上可逆 $n\times n$ 矩阵 $P$，使 $P^{-1}AP$ 处于有理标准形。这个有理标准形唯一。

\begin{proof}
令 $V=F^n$，让 $A$ 定义 $V$ 上的线性变换 $T(v)=Av$。由定理 14 可选 $V$ 的一组基，使 $T$ 的矩阵为有理标准形。若 $P$ 的列是这组新基向量在标准基下的坐标，则 $P^{-1}AP$ 就是该有理标准形。唯一性由不变因子的唯一性给出。
\end{proof}

\noindent\textbf{定义.} 域 $F$ 上 $n\times n$ 矩阵的不变因子，定义为由该矩阵定义的线性变换使 $F^n$ 成为 $F[x]$-模后所得的不变因子。

\noindent\textbf{定理 17.} 设 $A$ 与 $B$ 是域 $F$ 上的 $n\times n$ 矩阵。则 $A$ 与 $B$ 相似当且仅当它们有相同的不变因子。

\begin{proof}
这是定理 15 直接应用到 $V=F^n$ 的矩阵形式。相似矩阵给出同一个线性变换在不同基下的矩阵；而同一不变因子列表也由基本定理唯一地确定相应 $F[x]$-模的同构类，进而确定相似类。
\end{proof}

\noindent\textbf{推论 18.} 设 $A,B$ 是域 $F$ 上的两个 $n\times n$ 矩阵，并设 $F$ 是域 $K$ 的子域。
\begin{enumerate}[label=(\arabic*)]
\item 若 $A$ 与 $B$ 在 $F$ 上相似，则它们在 $K$ 上相似。
\item 若 $A$ 与 $B$ 在 $K$ 上相似，并且二者在 $K[x]$ 中的不变因子都事实上属于 $F[x]$，则 $A$ 与 $B$ 在 $F$ 上相似。
\end{enumerate}

\begin{proof}
(1) 若 $B=P^{-1}AP$ 且 $P$ 的条目在 $F$ 中，则把 $P$ 看作 $K$ 上矩阵同样给出 $K$ 上相似。 (2) 若在 $K$ 上相似，则由定理 17 它们在 $K[x]$ 中有相同的不变因子。若这些不变因子属于 $F[x]$，则它们也是作为 $F$ 上矩阵的不变因子；再由定理 17，$A$ 与 $B$ 在 $F$ 上相似。
\end{proof}

\noindent\textbf{引理 19.} 设 $a(x)\in F[x]$ 是任意首一多项式，$C(a)$ 是它的伴随矩阵。则 $C(a)$ 的最小多项式和特征多项式都等于 $a(x)$。更一般地，如果 $M$ 是块对角矩阵，对角块为 $A_1,\ldots,A_k$，则 $M$ 的特征多项式为这些 $A_i$ 的特征多项式之积，最小多项式为这些 $A_i$ 的最小多项式的最小公倍式。

\begin{proof}
伴随矩阵 $C(a)$ 表示 $F[x]/(a)$ 中“乘以 $x$”的线性变换；整个循环模的零化子正是 $(a)$，所以最小多项式为 $a$。其特征多项式可直接由 $xI-C(a)$ 的行列式展开得到，也可由后面的命题 20 得到。对于块对角矩阵，行列式按块相乘给出特征多项式的乘积；一个多项式零化整个块对角矩阵当且仅当它零化每个块，所以最小多项式是各块最小多项式的最小公倍式。
\end{proof}

\noindent\textbf{命题 20.} 设 $A$ 的不变因子为 $a_1(x),a_2(x),\ldots,a_m(x)$。则
\[
c_A(x)=a_1(x)a_2(x)\cdots a_m(x).
\]
特别地，$m_A(x)$ 整除 $c_A(x)$。

\begin{proof}
矩阵 $A$ 相似于有理标准形，其对角块是 $C(a_1),\ldots,C(a_m)$。相似矩阵有相同特征多项式；而块对角矩阵的特征多项式为各块特征多项式的乘积，且每个伴随矩阵 $C(a_i)$ 的特征多项式为 $a_i$。因此 $c_A=a_1\cdots a_m$。又最小多项式是最大不变因子 $a_m$，而 $a_m\mid a_1\cdots a_m$，故 $m_A\mid c_A$。
\end{proof}

下面给出计算矩阵不变因子的一种矩阵方法。命题 20 中的关系在确定矩阵 $A$ 的不变因子时经常很有用，特别是对低阶矩阵而言（参见练习 3、4 以及下面的例子）。以下结果依赖上一节练习 16--19；它在一般情形中计算不变因子，其证明在练习中概述。

设 $A$ 是域 $F$ 上的 $n\times n$ 矩阵。于是 $xI-A$ 是一个条目在 $F[x]$ 中的 $n\times n$ 矩阵。下列三种运算称为初等行、列运算：
\begin{enumerate}[label=(\alph*)]
\item 交换两行或两列；
\item 把一行或一列的一个 $F[x]$ 倍加到另一行或另一列上；
\item 把任一行或任一列乘以 $F[x]$ 中的一个单位，也就是乘以 $F$ 中一个非零元素。
\end{enumerate}

\noindent\textbf{定理 21.} 设 $A$ 是域 $F$ 上的 $n\times n$ 矩阵。利用上述三种初等行、列运算，条目在 $F[x]$ 中的矩阵 $xI-A$ 可以化为对角形（称为 $A$ 的 Smith 标准形）
\[
\begin{pmatrix}
1&&&&&\\
&\ddots&&&&\\
&&1&&&\\
&&&a_1(x)&&\\
&&&&\ddots&\\
&&&&&a_m(x)
\end{pmatrix},
\]
其中 $a_1(x),a_2(x),\ldots,a_m(x)$ 是 $F[x]$ 中首一非零元素，次数至少为 $1$，并满足
\[
a_1(x)\mid a_2(x)\mid\cdots\mid a_m(x).
\]
这些元素 $a_1(x),\ldots,a_m(x)$ 正是 $A$ 的不变因子。

\begin{proof}
参见练习。
\end{proof}

\paragraph{不变因子分解算法：化为有理标准形}
如上一节末尾的练习所述，跟踪把 $xI-A$ 对角化所需的运算，可以显式给出一个矩阵 $P$，使得 $P^{-1}AP$ 为有理标准形。等价地，若 $V$ 是一个给定的 $F[x]$-模，并且有向量空间基 $[e_1,e_2,\ldots,e_n]$，那么 $P$ 给出使 $V$ 分解为循环 $F[x]$-模直和的不变因子分解所需的换基。特别地，若 $A$ 是由 $x$ 定义的 $F[x]$-模 $V$ 的线性变换 $T$ 在这组基下的矩阵（即 $T(e_j)=xe_j=\sum_{i=1}^n a_{ij}e_i$，其中 $A=(a_{ij})$），那么矩阵 $P$ 给出 $V$ 的换基，使得 $T$ 的矩阵为有理标准形。

我们先在一般背景下描述算法：给定一个具有向量空间基 $[e_1,e_2,\ldots,e_n]$ 的 $F[x]$-模 $V$，求它的不变因子分解（证明在练习中概述）。然后描述怎样把给定的 $n\times n$ 矩阵 $A$ 化为有理标准形（这时不再显式提及底层向量空间和相应线性变换）。下面例 2 与例 3 给出该算法的具体数值例子。

\paragraph{不变因子分解算法}
设 $V$ 是具有向量空间基 $[e_1,e_2,\ldots,e_n]$ 的 $F[x]$-模（于是这些元素尤其是 $V$ 作为 $F[x]$-模的生成元）。设 $T$ 是由 $x$ 定义的 $V$ 到自身的线性变换，并设 $A$ 是 $T$ 关于 $V$ 的这组基的 $n\times n$ 矩阵，即
\[
T(e_j)=xe_j=\sum_{i=1}^n a_{ij}e_i,
\qquad A=(a_{ij}).
\]
\begin{enumerate}[label=(\arabic*)]
\item 使用以下三种初等行、列运算，在 $F[x]$ 上把矩阵 $xI-A$ 对角化，并记录所用的行运算：
\begin{enumerate}[label=(\alph*)]
\item 交换两行或两列（第 $i$ 行与第 $j$ 行交换记为 $R_i\leftrightarrow R_j$；列同理记为 $C_i\leftrightarrow C_j$）；
\item 把一行或一列的一个 $F[x]$ 倍加到另一行或列上（若把 $p(x)$ 倍第 $j$ 行加到第 $i$ 行，记为 $R_i+p(x)R_j\mapsto R_i$；列同理记为 $C_i+p(x)C_j\mapsto C_i$）；
\item 把任一行或任一列乘以 $F[x]$ 中的一个单位，也就是乘以 $F$ 中的一个非零元素（若把第 $i$ 行乘以 $u\in F^\times$，记为 $uR_i$；列同理记为 $uC_i$）。
\end{enumerate}
\item 从 $F[x]$-模生成元 $[e_1,e_2,\ldots,e_n]$ 出发，对 (1) 中使用的每一个行运算，按下列规则改变生成元集合：
\begin{enumerate}[label=(\alph*)]
\item 若第 $i$ 行与第 $j$ 行交换，则交换第 $i$ 个与第 $j$ 个生成元；
\item 若把 $p(x)$ 倍第 $j$ 行加到第 $i$ 行，则从第 $j$ 个生成元中减去 $p(x)$ 倍第 $i$ 个生成元（注意指标）；
\item 若第 $i$ 行乘以单位 $u\in F$，则把第 $i$ 个生成元除以 $u$。
\end{enumerate}
\item 当 $xI-A$ 已经被对角化为定理 21 中的形式时，$V$ 的生成元 $[e_1,e_2,\ldots,e_n]$ 已经成为原来 $e_1,e_2,\ldots,e_n$ 的 $F[x]$-线性组合。利用 $xe_j=T(e_j)=\sum_{i=1}^n a_{ij}e_i$，把这些元素写成 $e_1,e_2,\ldots,e_n$ 的 $F$-线性组合。对角化以后，前 $n-m$ 个这样的线性组合为 $0$（这给计算提供一个有用的数值检验），其余 $m$ 个线性组合非零，也就是说，$V$ 的生成元具有形式
\[
[0,\ldots,0,f_1,\ldots,f_m],
\]
并且它们精确对应于定理 21 中的对角元素。元素 $f_1,\ldots,f_m$ 是 $V$ 的不变因子分解中循环因子的 $F[x]$-模生成元（零化子分别为 $(a_1(x)),\ldots,(a_m(x))$）：
\[
V=F[x]f_1\oplus F[x]f_2\oplus\cdots\oplus F[x]f_m,
\qquad
F[x]f_i\cong F[x]/(a_i(x)),
\]
其中 $i=1,2,\ldots,m$。这给出 $F[x]$-模 $V$ 的不变因子分解。
\item 每个循环因子的相应向量空间基由元素
\[
f_i,Tf_i,T^2f_i,\ldots,T^{\deg a_i(x)-1}f_i
\]
给出。
\item 把 (4) 中算出的向量空间基的第 $k$ 个元素写成原向量空间基 $[e_1,e_2,\ldots,e_n]$ 的线性组合，并把所得坐标作为一个 $n\times n$ 矩阵 $P$ 的第 $k$ 列。则 $P^{-1}AP$ 处于有理标准形（对角块为多项式 $a_i(x)$ 的伴随矩阵）。这是线性变换 $T$ 关于 (4) 中向量空间基的矩阵。
\end{enumerate}

现在描述把给定 $n\times n$ 矩阵 $A$ 化为有理标准形的算法，即求 $n\times n$ 矩阵 $P$，使 $P^{-1}AP$ 处于有理标准形。这不过是把上面的算法应用到向量空间 $V=F^n$（$n\times1$ 列向量空间）及其标准基 $[e_1,e_2,\ldots,e_n]$，并令 $T$ 为由 $A$ 和这组基定义的线性变换。下面省略对底层向量空间与相应线性变换的显式提及，所以算法纯粹是矩阵论的。

\paragraph{把 $n\times n$ 矩阵化为有理标准形}
设 $A$ 是域 $F$ 上的 $n\times n$ 矩阵。
\begin{enumerate}[label=(\arabic*)]
\item 使用上述三种初等行、列运算，在 $F[x]$ 上把矩阵 $xI-A$ 对角化，并记录所用的行运算。定义 $d_1,\ldots,d_m$ 分别为对角线上出现的首一非常数多项式 $a_1(x),\ldots,a_m(x)$ 的次数。
\item 从 $n\times n$ 单位矩阵 $P'$ 出发，对 (1) 中使用的每一个行运算，按以下规则改变矩阵 $P'$：
\begin{enumerate}[label=(\alph*)]
\item 若 $R_i\leftrightarrow R_j$，则交换 $P'$ 的第 $i$ 列与第 $j$ 列（即对 $P'$ 作 $C_i\leftrightarrow C_j$）；
\item 若 $R_i+p(x)R_j\mapsto R_i$，则从 $P'$ 的第 $j$ 列中减去矩阵 $p(A)$ 乘以 $P'$ 的第 $i$ 列所得的列向量（即对 $P'$ 作 $C_j-p(A)C_i\mapsto C_j$；注意指标）；
\item 若 $uR_i$，则把 $P'$ 的第 $i$ 列的所有元素除以 $u$（即对 $P'$ 作 $u^{-1}C_i$）。
\end{enumerate}
\item 当 $xI-A$ 被对角化为定理 21 中的形式时，矩阵 $P'$ 的前 $n-m$ 列为 $0$（这给计算提供一个有用的数值检验），其余 $m$ 列非零。对每个 $i=1,2,\ldots,m$，把 $P'$ 的第 $i$ 个非零列向量依次乘以 $A^0=I,A^1,A^2,\ldots,A^{d_i-1}$，并按此顺序把所得列向量作为 $n\times n$ 矩阵 $P$ 接下来的 $d_i$ 列。则 $P^{-1}AP$ 处于有理标准形（其对角块为 (1) 中多项式 $a_1(x),\ldots,a_m(x)$ 的伴随矩阵）。
\end{enumerate}

在线性变换（或矩阵）的标准形理论中，特征多项式扮演有限阿贝尔群的阶的角色，而最小多项式扮演指数的角色（毕竟它们是同一类不变量：一个是主理想整环 $\mathbb Z$ 上模的不变量，另一个是主理想整环 $F[x]$ 上模的不变量），因此我们可以直接解决与第 5 章中有限阿贝尔群问题相对应的问题。特别包括：
\begin{enumerate}[label=(\Alph*)]
\item 确定给定矩阵的有理标准形（类似于把有限阿贝尔群分解为循环群直积）；
\item 判定两个给定矩阵是否相似（类似于判定两个给定有限阿贝尔群是否同构）；
\item 确定域 $F$ 上具有给定特征多项式的矩阵的所有相似类（类似于确定给定阶的所有阿贝尔群）；
\item 确定域 $F$ 上具有给定最小多项式的所有 $n\times n$ 矩阵相似类（类似于确定给定指数、秩至多为 $n$ 的所有阿贝尔群）。
\end{enumerate}

\paragraph{例}
\begin{enumerate}[label=(\arabic*)]
\item 我们求下列有理矩阵的有理标准形，并判定它们是否相似：
\[
A=\begin{pmatrix}2&-2&14\\0&3&-7\\0&0&2\end{pmatrix},\qquad
B=\begin{pmatrix}1&-44&-14\\0&85&28\\0&0&3\end{pmatrix},\qquad
C=\begin{pmatrix}1&2&2\\0&2&-1\\0&0&3\end{pmatrix}.
\]
直接计算表明三个矩阵有相同的特征多项式
\[
c_A(x)=c_B(x)=c_C(x)=(x-2)^2(x-3).
\]
由于最小多项式与特征多项式有相同根，最小多项式只可能是 $(x-2)(x-3)$ 或 $(x-2)^2(x-3)$。很快可算得
\[
(A-2I)(A-3I)=0,
\]
而 $(B-2I)(B-3I)\neq0$（其 $(1,1)$ 条目非零），且 $(C-2I)(C-3I)\neq0$（其 $(1,2)$ 条目非零）。因此
\[
m_A(x)=(x-2)(x-3),\qquad m_B(x)=m_C(x)=(x-2)^2(x-3).
\]
于是 $B$ 与 $C$ 不再有其他不变因子。由于 $A$ 的不变因子整除最小多项式且乘积为特征多项式，$A$ 的不变因子为
\[
x-2,\qquad (x-2)(x-3)=x^2-5x+6.
\]
（对 $2\times2$ 和 $3\times3$ 矩阵，特征多项式与最小多项式已足以确定所有不变因子，参见练习 3、4。）因此 $B$ 与 $C$ 相似，而二者都不与 $A$ 相似。有理标准形为（注意 $(x-2)^2(x-3)=x^3-7x^2+16x-12$）
\[
\begin{pmatrix}2&0&0\\0&0&-6\\0&1&5\end{pmatrix},
\qquad
\begin{pmatrix}0&0&12\\1&0&-16\\0&1&7\end{pmatrix},
\qquad
\begin{pmatrix}0&0&12\\1&0&-16\\0&1&7\end{pmatrix}.
\]

\item 上例的有理标准形只是通过确定矩阵的特征多项式与最小多项式得到的。如前所述，对 $2\times2$ 与 $3\times3$ 矩阵这已经足够，因为这些信息足以确定所有不变因子。然而对更大矩阵一般并不足够（参见下一例），还需要更多计算来确定不变因子。本例再次用上面两种算法计算例 (1) 中矩阵 $A$ 的有理标准形。对这个小矩阵来说这在计算上更麻烦（这一点会很明显），但即使在此情形它也有一个优点：它还显式算出矩阵 $P$，使 $P^{-1}AP$ 为有理标准形。

\noindent I.（不变因子分解）我们在 $\mathbb Q[x]$ 中用行、列运算把矩阵
\[
xI-A=\begin{pmatrix}
x-2&2&-14\\
0&x-3&7\\
0&0&x-2
\end{pmatrix}
\]
化为对角形。采用不变因子分解算法中的记号，$R_i\leftrightarrow R_j$ 表示交换第 $i$ 行与第 $j$ 行；若把 $a$ 倍第 $j$ 行加到第 $i$ 行，记作 $R_i+aR_j\mapsto R_i$；第 $i$ 行乘以 $u$ 简记作 $uR_i$（列运算同理，用 $C$ 代替 $R$）。在下面的化简中，最开始两个运算相当临时，只是为了使计算中各处都是整数。该化简把 $xI-A$ 化为对角矩阵，其非单位对角元为
\[
x-2,\qquad x^2-5x+6,
\]
因而确定了该矩阵的不变因子，这与例 (1) 所得一致。令 $V$ 为 $\mathbb Q$ 上三维向量空间，基为 $e_1,e_2,e_3$，并令 $T$ 为相应线性变换（它定义 $x$ 在 $V$ 上的作用），即
\[
xe_1=T(e_1)=2e_1,
\qquad
xe_2=T(e_2)=-2e_1+3e_2,
\qquad
xe_3=T(e_3)=14e_1-7e_2+2e_3.
\]
化简中用到的行运算为
\[
R_1+R_2\mapsto R_1,
\quad -R_1,
\quad R_2+(x-3)R_1\mapsto R_2,
\quad R_2-7R_3\mapsto R_2,
\quad R_2\leftrightarrow R_3.
\]
从 $V$ 的基 $[e_1,e_2,e_3]$ 出发，并按正文中的规则改变它，得到
\[
\begin{aligned}
[e_1,e_2,e_3]
&\mapsto [e_1,e_2-e_1,e_3]
\mapsto [-e_1,e_2-e_1,e_3]\\
&\mapsto [-e_1-(x-3)(e_2-e_1),e_2-e_1,e_3]\\
&\mapsto [-e_1-(x-3)(e_2-e_1),e_2-e_1,e_3+7(e_2-e_1)]\\
&\mapsto [-e_1-(x-3)(e_2-e_1),e_3+7(e_2-e_1),e_2-e_1].
\end{aligned}
\]
利用上面对 $x$ 作用的公式，可见这些最后的元素是
\[
[0,-7e_1+7e_2+e_3,-e_1+e_2],
\]
分别对应于 $xI-A$ 对角化形式中的元素 $1,x-2,x^2-5x+6$。因此
\[
f_1=-7e_1+7e_2+e_3,
\qquad
f_2=-e_1+e_2
\]
是 $V$ 作为 $\mathbb Q[x]$-模的不变因子分解中两个循环因子的 $\mathbb Q[x]$-模生成元。相应的 $\mathbb Q$-向量空间基为 $f_1$ 以及 $f_2,xf_2=Tf_2$，即
\[
-7e_1+7e_2+e_3,
\qquad
-e_1+e_2,
\qquad
T(-e_1+e_2)=-4e_1+3e_2.
\]
于是矩阵
\[
P=\begin{pmatrix}
-7&-1&-4\\
7&1&3\\
1&0&0
\end{pmatrix}
\]
把 $A$ 共轭到其有理标准形：
\[
P^{-1}AP=\begin{pmatrix}2&0&0\\0&0&-6\\0&1&5\end{pmatrix},
\]
这很容易直接检验。

\noindent II.（直接把 $A$ 化为有理标准形）利用对角化 $xI-A$ 时使用的行运算，按上面的算法确定矩阵 $P'$。这里 $d_1=1,d_2=2$，分别对应于 $P'$ 的第二与第三个非零列。因此 $P$ 的列分别为
\[
\begin{pmatrix}-7\\7\\1\end{pmatrix},
\qquad
\begin{pmatrix}-1\\1\\0\end{pmatrix},
\qquad
A\begin{pmatrix}-1\\1\\0\end{pmatrix}=\begin{pmatrix}-4\\3\\0\end{pmatrix},
\]
从而再次得到上面的矩阵 $P$。

\item 对 $3\times3$ 矩阵 $A$ 来说，仅为确定有理标准形（等价地，不变因子）而执行上面冗长计算并非必要，正如例 (1) 所示。然而对 $n\times n$ 矩阵，若 $n\geq4$，一般来说特征多项式与最小多项式的计算并不足以确定所有不变因子，因此前一例中更广泛的计算可能变得必要。比如考虑矩阵
\[
D=\begin{pmatrix}
1&0&0&-4\\
2&-1&0&-4\\
1&1&0&-2\\
0&0&1&3
\end{pmatrix}.
\]
简短计算表明 $D$ 的特征多项式为 $(x-1)^4$。可能的最小多项式为 $x-1,(x-1)^2,(x-1)^3,(x-1)^4$。显然 $D-I\neq0$，另一次简短计算给出 $(D-I)^2=0$，所以 $D$ 的最小多项式为 $(x-1)^2$。于是有两种可能的不变因子集合：
\[
x-1,x-1,(x-1)^2
\qquad\text{和}\qquad
(x-1)^2,(x-1)^2.
\]
为确定 $D$ 的不变因子，把前一例的方法应用于 $4\times4$ 矩阵
\[
xI-D=\begin{pmatrix}
x-1&0&0&4\\
-2&x+1&0&4\\
-1&-1&x&2\\
0&0&-1&x-3
\end{pmatrix}.
\]
由初等行、列运算从此矩阵得到的对角矩阵为
\[
\begin{pmatrix}
1&0&0&0\\
0&1&0&0\\
0&0&(x-1)^2&0\\
0&0&0&(x-1)^2
\end{pmatrix},
\]
这说明 $D$ 的不变因子为 $(x-1)^2,(x-1)^2$。一组把 $xI-D$ 对角化的初等行、列运算为
\[
\begin{gathered}
R_1\leftrightarrow R_3,
\quad -R_1,
\quad R_2+2R_1\mapsto R_2,
\quad R_3-(x-1)R_1\mapsto R_3,\\
C_3+(x-1)C_1\mapsto C_3,
\quad C_4+2C_1\mapsto C_4,
\quad R_2\leftrightarrow R_4,
\quad -R_2,\\
R_3+2R_2\mapsto R_3,
\quad R_4-(x+1)R_2\mapsto R_4,
\quad C_3+2C_2\mapsto C_3,
\quad C_4+(x-3)C_2\mapsto C_4.
\end{gathered}
\]
I.（不变因子分解）若在此情形中 $e_1,e_2,e_3,e_4$ 是 $V$ 的一组基，则像前一例一样使用这些行运算可见，对应于上述因子的 $V$ 的生成元为
\[
(x-1)e_1-2e_2-e_3=0,
\qquad
-2e_1+(x+1)e_2-e_4=0,
\qquad
e_1,
\qquad
e_2.
\]
因此在此情形中 $V$ 的不变因子分解中两个直因子的向量空间基为 $e_1,Te_1$ 以及 $e_2,Te_2$，其中 $T$ 是由 $D$ 定义的线性变换，即
\[
e_1,
\quad e_1+2e_2+e_3,
\quad e_2,
\quad 2e_1-e_2+e_4.
\]
相应的换基矩阵为
\[
P=\begin{pmatrix}
1&1&0&2\\
0&2&1&-1\\
0&1&0&0\\
0&0&0&1
\end{pmatrix},
\]
使得 $P^{-1}DP$ 为有理标准形，这也可直接检验。

II.（直接把 $D$ 化为有理标准形）与例 (2) 类似，由对角化 $xI-D$ 所用的行运算确定算法中的矩阵 $P'$。这里 $d_1=2,d_2=2$，分别对应于 $P'$ 的第三与第四个非零列。把这些列分别乘以 $I,D$，即可再次得到上面的矩阵 $P$。

\item 本例确定所有系数在 $\mathbb Q$ 中、特征多项式为 $(x^4-1)(x^2-1)$ 的矩阵 $A$ 的相似类。首先注意，具有次数 $6$ 的特征多项式的矩阵必为 $6\times6$ 矩阵。多项式 $(x^4-1)(x^2-1)$ 在 $\mathbb Q[x]$ 中分解为
\[
(x-1)^2(x+1)^2(x^2+1).
\]
由于 $A$ 的最小多项式 $m_A(x)$ 与 $c_A(x)$ 有相同根，可知 $(x-1)(x+1)(x^2+1)$ 整除 $m_A(x)$。设某个 $A$ 的不变因子为 $a_1(x),\ldots,a_m(x)$，于是 $a_m(x)=m_A(x)$，$a_i(x)\mid a_{i+1}(x)$（特别地，所有不变因子都整除 $m_A(x)$），且
\[
a_1(x)a_2(x)\cdots a_m(x)=(x^4-1)(x^2-1).
\]
容易看出，在这些限制下唯一允许的列表为
\begin{enumerate}[label=(\alph*)]
\item $(x-1)(x+1),\ (x-1)(x+1)(x^2+1)$；
\item $x-1,\ (x-1)(x+1)^2(x^2+1)$；
\item $x+1,\ (x-1)^2(x+1)(x^2+1)$；
\item $(x-1)^2(x+1)^2(x^2+1)$。
\end{enumerate}
现在可以容易写出相应的伴随矩阵直和，从而得到 $4$ 个相似类的代表。下一节将看到即使在 $M_6(\mathbb C)$ 中仍然只有 $4$ 个相似类。

\item 本例求所有满足 $A^6=I$ 且条目在 $\mathbb Q$ 中的 $3\times3$ 矩阵 $A$ 的相似类。对每个这样的 $A$，其最小多项式整除 $x^6-1$，而在 $\mathbb Q[x]$ 中该多项式的完全分解为
\[
x^6-1=(x-1)(x+1)(x^2-x+1)(x^2+x+1).
\]
反过来，若 $B$ 是任意 $3\times3$ 矩阵且其最小多项式整除 $x^6-1$，则 $B^6=I$。对 $B$ 的最小多项式的唯一限制是其次数至多为 $3$（由 Cayley--Hamilton 定理）。因此这种矩阵 $A$ 的最小多项式的唯一可能为
\begin{enumerate}[label=(\alph*)]
\item $x-1$；
\item $x+1$；
\item $x^2-x+1$；
\item $x^2+x+1$；
\item $(x-1)(x+1)$；
\item $(x-1)(x^2-x+1)$；
\item $(x-1)(x^2+x+1)$；
\item $(x+1)(x^2-x+1)$；
\item $(x+1)(x^2+x+1)$。
\end{enumerate}
在有理标准形的限制下，它们给出以下允许的不变因子列表：
\begin{enumerate}[label=(\roman*)]
\item $x-1,x-1,x-1$；
\item $x+1,x+1,x+1$；
\item $x-1,(x-1)(x+1)$；
\item $x+1,(x-1)(x+1)$；
\item $(x-1)(x^2-x+1)$；
\item $(x-1)(x^2+x+1)$；
\item $(x+1)(x^2-x+1)$；
\item $(x+1)(x^2+x+1)$。
\end{enumerate}
注意，如果最小多项式为 $x^2+x+1$ 或 $x^2-x+1$，则不存在合适的不变因子集合。现在可以写出相应的有理标准形；例如 (i) 是 $I$，(ii) 是 $-I$，而 (iii) 是
\[
\begin{pmatrix}
1&0&0\\
0&0&1\\
0&1&0
\end{pmatrix}.
\]
还要注意，这个结果也可表述为：任意条目在 $\mathbb Q$ 中且阶（当然是乘法阶）整除 $6$ 的 $3\times3$ 矩阵，都相似于这 $8$ 个矩阵之一；所以本例确定了群 $GL_3(\mathbb Q)$ 中阶为 $1,2,3,6$ 的所有元素（差一个相似）。
\end{enumerate}

\subsection*{练习}

\begin{enumerate}[label=\arabic*.]
\item 证明相似的线性变换（或 $n\times n$ 矩阵）有相同的特征多项式和相同的最小多项式。

\item 设 $M$ 如引理 19。证明 $M$ 的最小多项式是 $A_1,
\ldots,A_k$ 的最小多项式的最小公倍式。

\item 证明 $F$ 上两个不是标量矩阵的 $2\times2$ 矩阵相似当且仅当它们有相同的特征多项式。

\item 证明两个 $3\times3$ 矩阵相似当且仅当它们有相同的特征多项式和相同的最小多项式。给出一个 $4\times4$ 矩阵的明确反例，说明此命题对 $4\times4$ 矩阵不成立。

\item 直接利用 $n$ 维 $F$-向量空间 $V$ 到自身的全体线性变换构成维数为 $n^2$ 的 $F$-向量空间这一事实，证明线性变换 $T$ 的最小多项式次数至多为 $n^2$。

\item 证明 $n\times n$ 矩阵 $A$ 的特征多项式的常数项为 $(-1)^n\det A$，$x^{n-1}$ 的系数为 $A$ 的对角线元素之和的负数。对角线元素之和称为 $A$ 的迹，记为 $\operatorname{tr}(A)$。证明 $\det A$ 是 $A$ 的特征值之积，$\operatorname{tr}(A)$ 是 $A$ 的特征值之和。

\item 求矩阵
\[
\begin{pmatrix}
0&1&0&0\\
0&0&1&0\\
0&0&0&1\\
1&0&0&0
\end{pmatrix}
\]
的特征值。

\item 验证伴随矩阵
\[
\begin{pmatrix}
0&0&0&\cdots&0&-a_0\\
1&0&0&\cdots&0&-a_1\\
0&1&0&\cdots&0&-a_2\\
\vdots&\vdots&\vdots& &\vdots&\vdots\\
0&0&0&\cdots&1&-a_{n-1}
\end{pmatrix}
\]
的特征多项式为
\[
x^n+a_{n-1}x^{n-1}+\cdots+a_1x+a_0.
\]

\item 求下列矩阵的有理标准形：
\[
\begin{pmatrix}0&-1&-1\\0&0&0\\-1&0&0\end{pmatrix},\qquad
\begin{pmatrix}c&0&-1\\0&c&1\\-1&1&c\end{pmatrix},\qquad
\begin{pmatrix}
422&465&15&-30\\
-420&-463&-15&30\\
840&930&32&-60\\
-140&-155&-5&12
\end{pmatrix}.
\]

\item 求 $\mathbb Q$ 上所有最小多项式为 $(x+2)^2(x-1)$ 的 $6\times6$ 矩阵相似类。给出所有不变因子列表并写出若干相应矩阵即可。

\item 求 $\mathbb C$ 上所有特征多项式为 $(x^4-1)(x^2-1)$ 的 $6\times6$ 矩阵相似类。

\item 求 $\mathbb F_2$ 上所有满足 $A^6=I$ 的 $3\times3$ 矩阵 $A$ 的相似类（与正文中在 $\mathbb Q$ 上的答案比较）。同样，求所有满足 $B^{20}=I$ 的 $4\times4$ 矩阵 $B$ 的相似类。

\item 证明：在 $\mathbb Q[x]$ 中给定特征多项式时，$\mathbb Q$ 上 $3\times3$ 矩阵的相似类数等于任意 $\mathbb Q$ 的扩域上的相似类数。给出例子说明此结论对 $4\times4$ 矩阵一般不成立。

\item 求特征多项式为 $x^2(x^2+1)^2$ 的线性变换的所有可能有理标准形。

\item 在相似意义下求所有精确阶为 $4$ 的 $2\times2$ 有理矩阵（即 $M_2(\mathbb Q)$ 中的矩阵，阶指乘法阶）。若矩阵条目来自 $\mathbb C$，也做同样问题。

\item 证明在 $\mathbb F_{19}[x]$ 中
\[
x^5-1=(x-1)(x^2-4x+1)(x^2+5x+1).
\]
用此在相似意义下求所有 $\mathbb F_{19}$ 上乘法阶为 $5$ 的 $2\times2$ 矩阵。

\item 求 $GL_3(\mathbb F_2)$ 的共轭类代表元。提示：与第 6 章定理 15 和命题 14 比较。

\item 设 $V$ 是 $\mathbb Q$ 上有限维向量空间，$T$ 是 $V$ 上非奇异线性变换，并满足
\[
T^{-1}=T^2+T.
\]
证明 $\dim V$ 被 $3$ 整除。若 $\dim V=3$，证明所有这样的变换 $T$ 都相似。

\item 令 $V$ 为无限维实向量空间
\[
\mathbb R^\infty=\{(a_0,a_1,a_2,
\ldots)\mid a_0,a_1,a_2,
\ldots\in\mathbb R\}.
\]
定义 $T:V\to V$ 为
\[
T(a_0,a_1,a_2,
\ldots)=(0,a_0,a_1,a_2,
\ldots).
\]
证明 $T$ 没有特征向量。

\item 设 $\ell$ 为素数，
\[
\Phi_\ell(x)=\frac{x^\ell-1}{x-1}=x^{\ell-1}+x^{\ell-2}+\cdots+x+1\in\mathbb Z[x]
\]
为第 $\ell$ 个分圆多项式，它在 $\mathbb Q$ 上不可约（见推论 9.14 后的例 4）。本题确定 $\Phi_\ell(x)$ 模任意素数 $p$ 后因子的最小次数，特别确定它何时模 $p$ 不可约。（这实际上确定了 $\Phi_\ell(x)$ 模 $p$ 的完整分解，参见第 13.6 节练习 8。）
\begin{enumerate}[label=(\alph*)]
\item 证明若 $p=\ell$，则 $\Phi_\ell(x)$ 在 $\mathbb F_\ell[x]$ 中被 $x-1$ 整除。
\item 设 $p\ne\ell$，令 $f$ 为 $p$ 在 $\mathbb F_\ell^\times$ 中的阶，即 $f$ 是使 $p^f\equiv1\pmod\ell$ 的最小正整数。证明 $m=f$ 是群 $GL_m(\mathbb F_p)$ 含有阶为 $\ell$ 的元素 $A$ 的第一个 $m$ 值。提示：使用第 11.1 节末尾 $GL_m(\mathbb F_p)$ 阶数公式。
\item 证明 $\Phi_\ell(x)$ 在 $\mathbb F_p[x]$ 中不被任何次数小于 $f$ 的多项式整除。提示：考虑这种因子的伴随矩阵并使用 (b)。令 $m_A(x)\in\mathbb F_p[x]$ 为 (b) 中矩阵 $A$ 的最小多项式，推出 $m_A(x)$ 不可约、次数为 $f$，并整除 $\Phi_\ell(x)$。
\item 特别地，证明 $\Phi_\ell(x)$ 模 $p$ 不可约当且仅当 $\ell-1$ 是使 $p$ 模 $\ell$ 同余于 $1$ 的最小幂，即 $p$ 是模 $\ell$ 的原根。
\end{enumerate}

\item 证明定理 21 前描述的前两种初等行、列操作不改变矩阵行列式，第三种初等操作只使行列式乘以一个单位。由定理 21 推出 $A$ 的特征多项式与 $A$ 的不变因子乘积只差一个单位。由于二者按定义都是首一多项式，推出它们相等（这给出命题 20 的另一证明）。

\medskip
以下练习勾勒定理 21 的证明。它们在 Euclidean 整环 $F[x]$ 上显式执行上一节练习 16--19 中描述的构造。设 $V$ 为 $n$ 维向量空间，基为 $v_1,v_2,\ldots,v_n$，$T$ 是由矩阵 $A$ 和此基定义的线性变换，即
\[
T(v_j)=\sum_{i=1}^n a_{ij}v_i,
\qquad j=1,2,
\ldots,n,
\]
其中 $A=(a_{ij})$。令 $F[x]^n$ 为秩 $n$ 的自由 $F[x]$-模，基为 $\xi_1,\xi_2,
\ldots,\xi_n$。有自然满 $F[x]$-模同态
\[
\varphi:F[x]^n\to V,
\qquad \xi_i\mapsto v_i.
\]
由于 $x$ 在 $V$ 上按 $T$ 作用，故
\[
x(v_j)=\sum_{i=1}^n a_{ij}v_i.
\]

\item 证明元素
\[
\nu_j=-a_{1j}\xi_1-\cdots-a_{j-1,j}\xi_{j-1}+(x-a_{jj})\xi_j-a_{j+1,j}\xi_{j+1}-\cdots-a_{nj}\xi_n
\]
对 $j=1,2,
\ldots,n$ 都属于 $\ker\varphi$。

\item
\begin{enumerate}[label=(\alph*)]
\item 证明 $x\xi_j=\nu_j+f_j$，其中 $f_j\in F\xi_1+\cdots+F\xi_n$ 是由 $\xi_1,
\ldots,
\xi_n$ 张成的 $F$-向量空间中的元素。
\item 证明
\[
F[x]\xi_1+\cdots+F[x]\xi_n
=(F[x]\nu_1+\cdots+F[x]\nu_n)+(F\xi_1+\cdots+F\xi_n).
\]
\end{enumerate}

\item 证明 $\nu_1,
\nu_2,
\ldots,
\nu_n$ 生成 $\ker\varphi$。提示：用上一题说明 $\ker\varphi$ 的任意元素可写成由 $\nu_1,
\ldots,
\nu_n$ 生成的模中的元素加上形如 $b_1\xi_1+\cdots+b_n\xi_n$ 的元素，其中 $b_i\in F$；再证明这种元素在 $\ker\varphi$ 中当且仅当所有 $b_i=0$，因为 $v_1,
\ldots,v_n$ 是 $V$ 的 $F$-基。

\item 证明 $\ker\varphi$ 的生成元 $\nu_1,
\ldots,
\nu_n$ 对应的关系矩阵为
\[
\begin{pmatrix}
 x-a_{11}&-a_{21}&\cdots&-a_{n1}\\
 -a_{12}&x-a_{22}&\cdots&-a_{n2}\\
 \vdots&\vdots&\ddots&\vdots\\
 -a_{1n}&-a_{2n}&\cdots&x-a_{nn}
\end{pmatrix}=xI-A^t.
\]
推出定理 21 以及求 $A$ 的不变因子的算法可由上一节练习 16--19 得到（注意，对此关系矩阵的行、列操作正是对定理 21 中矩阵作列、行操作，这解释了不变因子算法为什么要跟踪所用的行操作）。
\end{enumerate}



\section{Jordan 标准形}

继续沿用上一节记号：$F$ 是域，$F[x]$ 是 $F$ 上多项式环，$V$ 是 $F$ 上维数为 $n$ 的有限维向量空间，$T$ 是固定线性变换，由此把 $V$ 变为 $F[x]$-模；$A$ 是 $F$ 上的 $n\times n$ 矩阵。回忆，一旦固定 $V$ 的一组基，任意线性变换 $T$ 都定义一个矩阵 $A$；反之任意矩阵 $A$ 都定义一个线性变换 $T$。

上一节使用 P.I.D. $F[x]$ 上有限生成模基本定理的不变因子形式，得到线性变换和矩阵的有理标准形。本节使用基本定理的初等因子形式得到 Jordan 标准形。我们将看到，处于这个标准形的矩阵尽可能接近对角矩阵，因此它们比有理标准形中的矩阵更简单（但会失去某些“有理性”结果）。

模的初等因子是不变因子的素幂因子（推论 10）。对 $F[x]$-模 $V$ 来说，不变因子是次数至少为 $1$ 的首一多项式 $a_1(x),a_2(x),\ldots,a_m(x)$，并满足 $a_1\mid a_2\mid\cdots\mid a_m$；相应初等因子就是这些多项式的不可约因子的幂。和不变因子一样，我们要求它们为首一，从而唯一。

为了得到最简单的初等因子，我们假设多项式 $a_1(x),a_2(x),\ldots,a_m(x)$ 全部分解为一次因子，也就是说，$V$ 的初等因子都是形如 $(x-\lambda)^k$ 的线性多项式的幂。由于初等因子的乘积是特征多项式，这等价于假设域 $F$ 含有线性变换 $T$（或表示它的矩阵 $A$）的所有特征值。

在这个对 $F$ 的假设下，由定理 6 立即可知，$V$ 是有限多个循环 $F[x]$-模
\[
F[x]/(x-\lambda)^k
\]
的直和，其中 $\lambda\in F$ 是 $T$ 的某个特征值，这些直和因子对应于 $V$ 的初等因子。

现在为 $V$ 的每个初等因子直和因子选择一组向量空间基，使得 $T$ 的相应矩阵特别简单。回忆，由 $F[x]$-模结构的定义，$V$ 上的线性变换 $T$ 就是每个直和因子 $F[x]/(x-\lambda)^k$ 上元素 $x$ 的乘法作用。

考虑商模 $F[x]/(x-\lambda)^k$ 中的元素
\[
(\bar x-\lambda)^{k-1},(\bar x-\lambda)^{k-2},\ldots,\bar x-\lambda,1.
\]
把这些多项式展开为 $x$ 的多项式，可见它们与 $F$-基 $\bar x^{k-1},\bar x^{k-2},\ldots,\bar x,1$ 之间的转换矩阵是一个对角元全为 $1$ 的上三角矩阵。因此这个矩阵可逆（行列式为 $1$），所列元素确为 $F[x]/(x-\lambda)^k$ 的一组 $F$-基。关于这组基，乘以 $x$ 的线性变换按特别简单的方式作用（注意 $x=\lambda+(x-\lambda)$，且 $(\bar x-\lambda)^k=0$）：
\[
(\bar x-\lambda)^{k-1}\mapsto \lambda(\bar x-\lambda)^{k-1},
\]
\[
(\bar x-\lambda)^{k-2}\mapsto \lambda(\bar x-\lambda)^{k-2}+(\bar x-\lambda)^{k-1},
\]
依此类推，最后
\[
1\mapsto \lambda\cdot1+(\bar x-\lambda).
\]
因此关于这组基，乘以 $x$ 的矩阵为
\[
J_k(\lambda)=
\begin{pmatrix}
\lambda&1&0&\cdots&0\\
0&\lambda&1&\cdots&0\\
\vdots&\vdots&\ddots&\ddots&\vdots\\
0&0&\cdots&\lambda&1\\
0&0&\cdots&0&\lambda
\end{pmatrix}.
\]

\noindent\textbf{定义.} 上述 $k\times k$ 矩阵称为特征值为 $\lambda$ 的 $k\times k$ \emph{初等 Jordan 矩阵}，或称特征值为 $\lambda$、大小为 $k$ 的 \emph{Jordan 块}。

把这应用到 $V$ 的初等因子分解中的每个循环因子上，就得到 $V$ 的一组向量空间基，使得 $T$ 的矩阵是对应于 $V$ 初等因子的 Jordan 块的直和，即对角线上是 Jordan 块的块对角矩阵。这个矩阵由 $F[x]$-模 $V$ 的初等因子唯一决定（至多差一个对角块的排列）；反过来，由定理 9，初等因子列表也唯一决定 $V$ 的 $F[x]$-模同构类。

\noindent\textbf{定义.}
\begin{enumerate}[label=(\arabic*)]
\item 若一个矩阵是块对角矩阵，并且对角线上各块均为 Jordan 块，则称该矩阵处于\emph{Jordan 标准形}。
\item 线性变换 $T$ 的一个 \emph{Jordan 标准形}是一个表示 $T$ 且处于 Jordan 标准形的矩阵。
\end{enumerate}

上面的讨论已经证明了任意满足条件的线性变换 $T$ 都有 Jordan 标准形。与有理标准形情形一样，由初等因子的唯一性可知，Jordan 标准形在对角线上 Jordan 块排列次序以外唯一。因此称它为 $T$ 的 Jordan 标准形。总结如下：

\noindent\textbf{定理 22（线性变换的 Jordan 标准形）.} 设 $V$ 是域 $F$ 上有限维向量空间，$T$ 是 $V$ 上线性变换。假设 $F$ 含有 $T$ 的所有特征值。
\begin{enumerate}[label=(\arabic*)]
\item 存在 $V$ 的一组基，使得 $T$ 的矩阵处于 Jordan 标准形；也就是说，它是块对角矩阵，对角块是 $V$ 的初等因子所对应的 Jordan 块。
\item $T$ 的 Jordan 标准形在对角线上 Jordan 块的排列次序以外唯一。
\end{enumerate}

相应的矩阵形式如下：

\noindent\textbf{定理 23（矩阵的 Jordan 标准形）.} 设 $A$ 是域 $F$ 上 $n\times n$ 矩阵，并假设 $F$ 含有 $A$ 的所有特征值。
\begin{enumerate}[label=(\arabic*)]
\item 矩阵 $A$ 相似于一个处于 Jordan 标准形的矩阵；也就是说，存在 $F$ 上可逆 $n\times n$ 矩阵 $P$，使得 $P^{-1}AP$ 是块对角矩阵，其对角块是 $A$ 的初等因子对应的 Jordan 块。
\item $A$ 的 Jordan 标准形在对角线上 Jordan 块的排列次序以外唯一。
\end{enumerate}

Jordan 标准形与对角矩阵的差异仅在于第一超对角线上可能出现一些 $1$（且只有当存在大小大于 $1$ 的 Jordan 块时才会出现）。因此它接近对角矩阵。下面结果尤其说明，矩阵的 Jordan 标准形确实是相似意义下尽可能接近对角矩阵的形式。

\noindent\textbf{推论 24.}
\begin{enumerate}[label=(\arabic*)]
\item 如果矩阵 $A$ 相似于对角矩阵 $D$，则 $D$ 就是 $A$ 的 Jordan 标准形。
\item 两个对角矩阵相似当且仅当它们的对角元只差一个排列。
\end{enumerate}

\begin{proof}
第一断言由 Jordan 标准形唯一性立即得到，因为对角矩阵本身已经处于 Jordan 标准形（所有 Jordan 块都为大小 $1$）。Jordan 标准形的唯一性也给出 (2)。
\end{proof}

下一推论给出判断矩阵能否对角化的准则。

\noindent\textbf{推论 25.} 若 $A$ 是 $F$ 上 $n\times n$ 矩阵，且 $F$ 含有 $A$ 的所有特征值，则 $A$ 在 $F$ 上相似于对角矩阵，当且仅当 $A$ 的最小多项式没有重根。

\begin{proof}
若 $A$ 相似于对角矩阵，则对角矩阵的最小多项式没有重根（它的根正是对角线上出现的不同元素）。相似矩阵有相同最小多项式，所以 $A$ 的最小多项式没有重根。

反过来，假设 $A$ 的最小多项式没有重根，并令 $B$ 是 $A$ 的 Jordan 标准形。矩阵 $B$ 是以初等 Jordan 矩阵为对角块的块对角矩阵。由上一节末尾的练习，$B$ 的最小多项式是各 Jordan 块最小多项式的最小公倍式。大小为 $k$、特征值为 $\lambda$ 的 Jordan 块的最小多项式是 $(x-\lambda)^k$（这也可直接从该 Jordan 块给出一个零化子为 $((x-\lambda)^k)$ 的循环 $F[x]$-子模看出）。因此如果最小多项式无重根，所有 Jordan 块的大小都必须为 $1$。于是 $B$ 是对角矩阵，故 $A$ 可对角化。
\end{proof}

有理标准形与 Jordan 标准形之间的转换，就是第 5.2 节有限阿贝尔群情形中从不变因子到初等因子以及反向的同一算法。简言之，初等因子是不变因子的素幂因子；把每个不变因子分解为不同一次因子的幂，所得线性多项式幂的集合就是初等因子。反过来，给定初等因子列表时，先按不同特征值分成若干列，在每列内按次数非降排列；再给较短的列前面补常数多项式 $1$，使所有列长度相同；然后把各列第 $i$ 个多项式相乘，得到第 $i$ 个不变因子。

\paragraph{初等因子分解算法：化为 Jordan 标准形}
定理 21 给出计算任意矩阵 $A$ 不变因子的过程。把这些不变因子分解为线性因子的幂便得到 $A$ 的初等因子，从而决定 $A$ 的 Jordan 标准形。

更具体地说，不变因子分解算法从 $V$ 的基 $e_1,\ldots,e_n$ 出发，产生 $V$ 的若干元素 $f_1,\ldots,f_m$，它们是 $V$ 的不变因子分解中循环 $F[x]$-因子的 $F[x]$-模生成元，对应的零化子为 $(a_1(x)),\ldots,(a_m(x))$。由于初等因子分解由对每个循环模 $F[x]/(a_i(x))$ 应用中国剩余定理得到，这也给出初等因子分解中各循环因子的生成元。若
\[
a(x)=(x-\lambda_1)^{\alpha_1}(x-\lambda_2)^{\alpha_2}\cdots(x-\lambda_s)^{\alpha_s}
\]
且 $f$ 是对应于不变因子 $a(x)$ 的循环因子生成元，则
\[
\frac{a(x)}{(x-\lambda_1)^{\alpha_1}}f,
\ldots,
\frac{a(x)}{(x-\lambda_s)^{\alpha_s}}f
\]
分别给出对应初等因子 $(x-\lambda_i)^{\alpha_i}$ 的循环因子生成元。若
\[
g_i=\frac{a(x)}{(x-\lambda_i)^{\alpha_i}}f,
\]
则对应循环因子的向量空间基为
\[
(T-\lambda_i)^{\alpha_i-1}g_i,
(T-\lambda_i)^{\alpha_i-2}g_i,
\ldots,
(T-\lambda_i)g_i,g_i.
\]
把这些向量在原基下的坐标作为矩阵 $P$ 的列，则 $P^{-1}AP$ 为 Jordan 标准形。

\paragraph{例}
设 $A,B,C$ 是上一节例 1 中的矩阵，并取 $F=\mathbb Q$。注意 $\mathbb Q$ 含有这些矩阵的所有特征值。由于已经求出不变因子，可立即得到初等因子：$A$ 的初等因子为 $x-2,x-2,x-3$；$B$ 与 $C$ 的初等因子为 $(x-2)^2$ 和 $x-3$。因此 $A$ 相似于对角矩阵，而 $B,C$ 不可对角化。

\paragraph{例}
对上一节矩阵 $A$ 的显式基，已得到不变因子分解中的生成元
\[
f_1=-7e_1+7e_2+e_3,
\qquad f_2=-e_1+e_2,
\]
对应不变因子 $x-2$ 与 $(x-2)(x-3)$。由上述算法，初等因子分解中的生成元为
\[
f_1,\quad (x-3)f_2,\quad (x-2)f_2,
\]
分别对应 $x-2,x-2,x-3$。直接计算得这些向量为 $-7e_1+7e_2+e_3$、$-e_1$、$-2e_1+e_2$。以它们为列组成矩阵 $P$，则 $P^{-1}AP$ 为相应 Jordan 标准形。

\paragraph{例}
若一个 $4\times4$ 矩阵的不变因子为 $(x-1)^2,(x-1)^2$，相应 $\mathbb Q[x]$-模生成元为 $f_1=e_1$ 与 $f_2=e_2$，则这些不变因子本身就是初等因子。相应 Jordan 基为 $(T-I)f_1,f_1,(T-I)f_2,f_2$；在原书给定的具体矩阵中，这些向量为 $2e_2+e_3,e_1,2e_1-e_2+e_4,e_2$，从而给出两个大小为 $2$、特征值为 $1$ 的 Jordan 块。

当扩张底域时，相似类的数目可能增加（由推论 18(2)，永远不会减少）。例如在 $\mathbb C$ 上，特征多项式 $(x^4-1)(x^2-1)$ 的矩阵可进一步按一次因子 $x-i$ 与 $x+i$ 分解 Jordan 标准形；而在 $\mathbb Q$ 上这些因子合并为不可约二次因子 $x^2+1$，只能在有理标准形中体现。
\subsection*{练习}

\begin{enumerate}[label=\arabic*.]
\item 设向量空间 $V$ 是若干循环 $F[x]$-模的直和，其零化子分别为
\[
(x+1)^2,
\qquad (x-1)(x^2+1)^2,
\qquad (x^4-1),
\qquad (x+1)(x^2-1).
\]
求 $V$ 的不变因子和初等因子。

\item 证明：若 $\lambda_1,
\ldots,\lambda_n$ 是 $n\times n$ 矩阵 $A$ 的特征值，则对任意 $k\geq0$，$\lambda_1^k,
\ldots,
\lambda_n^k$ 是 $A^k$ 的特征值。

\item 用上面例 2 的方法，求显式矩阵 $P_1,P_2$，使 $P_1^{-1}BP_1$ 与 $P_2^{-1}CP_2$ 为 Jordan 标准形。用此显式构造一个矩阵 $Q$，它把 $B$ 共轭为 $C$，从而直接证明这两个矩阵相似。

\item 证明矩阵
\[
\begin{pmatrix}
9&4&5\\
-4&0&-3\\
-6&-4&-2
\end{pmatrix}
\]
的 Jordan 标准形正是本章开头给出的形式。显式求出一个矩阵 $P$，使它把该矩阵共轭到其 Jordan 标准形。说明为什么该矩阵不能对角化。

\item 求矩阵
\[
\begin{pmatrix}
1&0&0\\
0&0&-2\\
0&1&3
\end{pmatrix}
\]
的 Jordan 标准形。

\item 判断下列矩阵中哪些相似：
\[
\begin{pmatrix}-1&4&-4\\2&-1&3\\0&-4&3\end{pmatrix},\quad
\begin{pmatrix}-3&-4&0\\2&3&0\\8&8&1\end{pmatrix},\quad
\begin{pmatrix}-3&2&-4\\2&1&0\\3&-1&3\end{pmatrix},\quad
\begin{pmatrix}-1&4&-4\\0&-3&2\\0&-4&3\end{pmatrix}.
\]

\item 求下列矩阵的 Jordan 标准形：
\[
\begin{pmatrix}5&4&1\\-1&0&0\\-3&-4&1\end{pmatrix},\qquad
\begin{pmatrix}3&4&2\\-2&-3&-1\\-4&-4&-3\end{pmatrix}.
\]

\item 证明矩阵
\[
A=\begin{pmatrix}5&6&0\\-3&-4&0\\-2&0&1\end{pmatrix},
\qquad
B=\begin{pmatrix}3&-1&2\\-10&6&-14\\-6&3&-7\end{pmatrix}
\]
相似。证明 $A$ 与 $B$ 都可对角化，并求显式矩阵 $P_1,P_2$，使 $P_1^{-1}AP_1$ 与 $P_2^{-1}BP_2$ 为对角形式。

\item 证明矩阵
\[
A=\begin{pmatrix}-8&-10&-1\\7&9&1\\3&2&0\end{pmatrix},
\qquad
B=\begin{pmatrix}-3&2&-4\\4&-1&4\\4&-2&5\end{pmatrix}
\]
都以 $(x-1)^2(x+1)$ 为特征多项式，但其中一个可对角化而另一个不可对角化。求两个矩阵的 Jordan 标准形。

\item 求复数域上所有 $2\times2$、$3\times3$ 与 $4\times4$ 矩阵的 Jordan 标准形。

\item 验证矩阵
\[
A=\begin{pmatrix}
1&0&0&0\\
0&1&0&0\\
-2&-2&0&1\\
-2&0&-1&-2
\end{pmatrix}
\]
的特征多项式在 $\mathbb Q$ 上分解为一次因子。求 $A$ 在 $\mathbb Q$ 上的有理标准形与 Jordan 标准形。

\item 求矩阵
\[
\begin{pmatrix}
1&2&0&0\\
0&1&2&0\\
0&0&1&2\\
0&0&0&1
\end{pmatrix}
\]
的 Jordan 标准形。

\item 求矩阵
\[
\begin{pmatrix}
3&0&-2&-3\\
4&-8&14&-15\\
2&-4&7&-7\\
0&2&-4&3
\end{pmatrix}
\]
的 Jordan 标准形。

\item 证明矩阵
\[
A=\begin{pmatrix}
2&0&0&0\\
-4&-1&-4&0\\
2&1&3&0\\
-2&4&9&1
\end{pmatrix},
\qquad
B=\begin{pmatrix}
5&0&-4&-7\\
3&-8&15&-13\\
2&-4&7&-7\\
1&2&-5&1
\end{pmatrix}
\]
相似。

\item 证明矩阵
\[
A=\begin{pmatrix}
0&1&1&1\\
1&0&1&1\\
1&1&0&1\\
1&1&1&0
\end{pmatrix},
\qquad
B=\begin{pmatrix}
5&2&-8&-8\\
-6&-3&8&8\\
-3&-1&3&4\\
3&1&-4&-5
\end{pmatrix}
\]
都以 $(x-3)(x+1)^3$ 为特征多项式。判断它们是否相似，并求每个矩阵的 Jordan 标准形。

\item 求矩阵
\[
\begin{pmatrix}
1&1&1&1\\
0&1&0&-1\\
0&0&1&1\\
0&0&0&1
\end{pmatrix}
\]
的 Jordan 标准形，并求一个矩阵 $P$，使它把该矩阵共轭到其 Jordan 标准形。

\item 证明任意矩阵 $A$ 都与其转置 $A^t$ 相似。

\item 求特征多项式为 $(x-2)^3(x-3)^2$ 的线性变换的所有可能 Jordan 标准形。

\item 证明：所有特征多项式为 $f(x)$ 的 $n\times n$ 矩阵彼此相似，当且仅当 $f(x)$ 在 $F[x]$ 中的唯一因子分解没有重因子。

\item 证明下列两个 $M_p(\mathbb F_p)$ 中的矩阵相似：
\[
\begin{pmatrix}
0&0&0&\cdots&0&1\\
1&0&0&\cdots&0&0\\
0&1&0&\cdots&0&0\\
\vdots&\vdots&\vdots& &\vdots&\vdots\\
0&0&0&\cdots&1&0
\end{pmatrix}
\quad\text{和}\quad
\begin{pmatrix}
1&1&0&\cdots&0&0\\
0&1&1&\cdots&0&0\\
0&0&1&\cdots&0&0\\
\vdots&\vdots&\vdots& &\vdots&\vdots\\
0&0&0&\cdots&1&1\\
0&0&0&\cdots&0&1
\end{pmatrix}.
\]

\item 证明若 $A^2=A$，则 $A$ 相似于一个只以 $0$ 和 $1$ 为对角元的对角矩阵。

\item 证明复数域上满足 $A^3=A$ 的 $n\times n$ 矩阵 $A$ 可对角化。这个结论在任意域 $F$ 上是否成立？

\item 设 $A$ 是 $\mathbb Q$ 上 $2\times2$ 矩阵，满足 $A^3=I$ 但 $A\ne I$。把 $A$ 写成有理标准形；再把它视为 $\mathbb C$ 上矩阵，写出 Jordan 标准形。

\item 证明不存在 $\mathbb Q$ 上 $3\times3$ 矩阵 $A$ 满足 $A^8=I$ 但 $A^4\ne I$。

\item 求 $\mathbb Q$ 上所有条目都等于 $1$ 的 $n\times n$ 矩阵的 Jordan 标准形。

\item 求 $\mathbb F_p$ 上所有条目都等于 $1$ 的 $n\times n$ 矩阵的 Jordan 标准形（答案取决于 $p$ 是否整除 $n$）。

\item 求 $\mathbb Q$ 上 $n\times n$ 矩阵的 Jordan 标准形，其中所有非对角线条目均为 $1$，而主对角线条目均为 $0$。

\item 求 $\mathbb F_p$ 上 $n\times n$ 矩阵的 Jordan 标准形，其中所有非对角线条目均为 $1$，而主对角线条目均为 $0$。

与同一特征值 $\lambda$ 对应的、所有初等因子为 $(x-\lambda)$ 的幂的循环子模的直和，称为 $T$ 关于特征值 $\lambda$ 的\emph{广义特征空间}。这是 $F[x]$ 中素元 $p=x-\lambda$ 的 primary 分量，包含所有被某个 $(T-\lambda)$ 的幂零化的 $V$ 中元素。

\item 设 $V_i$ 是 $T$ 对特征值 $\lambda_i$ 的广义特征空间。对任意 $k\geq0$，证明 $T-\lambda_i$ 在子空间 $(T-\lambda_i)^kV_i$ 上的零化度，等于 $T-\lambda_i$ 在 $(T-\lambda_i)^kV$ 上的零化度，并且等于 $T$ 中特征值为 $\lambda_i$、大小大于 $k$ 的 Jordan 块个数（当 $k=0$ 时给出 Jordan 块的个数）。

\item 设 $\lambda$ 是域 $F$ 上有限维向量空间 $V$ 上线性变换 $T$ 的特征值。令
\[
r_k=\dim_F (T-\lambda)^kV
\]
为 $(T-\lambda)^k$ 在 $V$ 上的秩。证明对任意 $k\geq1$，
\[
r_{k-1}-2r_k+r_{k+1}
\]
是 $T$ 中对应于 $\lambda$ 且大小为 $k$ 的 Jordan 块个数（使用第 1 节练习 12）。这给出一种通过计算矩阵 $(A-\lambda I)^k$ 的秩来确定 Jordan 标准形的有效方法。

\item 设 $N$ 是域 $F$ 上 $n\times n$ 矩阵。若某个幂 $N^k$ 为零矩阵，则称 $N$ 为\emph{幂零矩阵}。证明任意幂零矩阵都相似于一个块对角矩阵，其每个块在第一超对角线上全为 $1$，其余条目全为 $0$。

\item 证明若 $N$ 是 $n\times n$ 幂零矩阵，则事实上 $N^n=0$。

\item 设 $A$ 是严格上三角的 $n\times n$ 矩阵（主对角线及其下方条目均为 $0$）。证明 $A$ 是幂零矩阵。

\item 证明 $n\times n$ 幂零矩阵的迹为 $0$。

\item 对 $0\leq i\leq n$，令 $d_i$ 为 $xI-A$ 的所有 $i\times i$ 子式行列式的最大公因子（把 $0\times0$ 子式定义为 $1$）。证明 $A$ 的 Smith 标准形对角线上第 $i$ 个元素为 $d_i/d_{i-1}$。这给出 $A$ 的不变因子。提示：证明这些最大公因子在初等行、列操作下不变。

\item 令 $V=\mathbb C^n$ 为通常的 $n$ 维复向量空间，$T$ 为线性变换
\[
T(a_1,a_2,
\ldots,a_n)=(0,a_1,a_2,
\ldots,a_{n-1}).
\]
求 $T$ 的 Jordan 标准形。

\item 设 $J$ 是复数域上大小为 $n$、特征值为 $\lambda$ 的 Jordan 块。
\begin{enumerate}[label=(\alph*)]
\item 证明若 $\lambda\ne0$，则 $J^2$ 的 Jordan 标准形是大小为 $n$、特征值为 $\lambda^2$ 的 Jordan 块。
\item 若 $\lambda=0$，证明 $J^2$ 的 Jordan 标准形有两个块，特征值均为 $0$；当 $n$ 为偶数时，两个块大小为 $n/2,n/2$，当 $n$ 为奇数时，大小为 $(n-1)/2,(n+1)/2$。
\end{enumerate}

\item 求矩阵 $A\in M_n(\mathbb C)$ 有平方根的充要条件，即存在 $B\in M_n(\mathbb C)$ 使 $A=B^2$。提示：设 $B$ 为 Jordan 标准形，并使用上一题考虑 $B^2$ 的 Jordan 标准形。

\item 设 $J$ 是特征为 $2$ 的域 $F$ 上大小为 $n$、特征值为 $\lambda$ 的 Jordan 块。求 $J^2$ 的 Jordan 标准形。求 $A\in M_n(F)$ 有平方根的充要条件，即存在 $B\in M_n(F)$ 使 $A=B^2$。

\medskip
以下练习考察矩阵函数（幂级数）并介绍 Jordan 标准形在常微分方程理论中的一些应用。以下矩阵均为 $K$ 上 $n\times n$ 矩阵，其中 $K=\mathbb R$ 或 $\mathbb C$。令
\[
G(x)=\sum_{k=0}^{\infty}\alpha_kx^k
\]
为 $K$ 系数幂级数，令 $G_N(x)=\sum_{k=0}^N\alpha_kx^k$。对 $A\in M_n(K)$，把 $A$ 代入 $G_N$ 得 $G_N(A)$。若对每个 $i,j$，$G_N(A)$ 的 $(i,j)$ 条目收敛到矩阵 $C$ 的 $(i,j)$ 条目，则称
\[
G(A)=\sum_{k=0}^{\infty}\alpha_kA^k
\]
收敛到 $C$，记作 $G(A)=C$。定义
\[
\|A\|=\sum_{i,j=1}^n |a_{ij}|.
\]

\item 证明对所有 $A,B\in M_n(K)$ 和 $\alpha\in K$，有：
\begin{enumerate}[label=(\alph*)]
\item $\|A+B\|\leq \|A\|+\|B\|$；
\item $\|AB\|\leq \|A\|\,\|B\|$；
\item $\|\alpha A\|=|\alpha|\,\|A\|$。
\end{enumerate}

\item 设 $R$ 为实或复幂级数 $G(x)$ 的收敛半径（若 $G(x)$ 对所有 $x\in K$ 收敛，则 $R=\infty$）。
\begin{enumerate}[label=(\alph*)]
\item 证明若 $\|A\|<R$，则 $G(A)$ 收敛。
\item 推出对所有矩阵 $A$，下列幂级数都收敛：
\[
\sin(A)=A-\frac{A^3}{3!}+\frac{A^5}{5!}+\cdots+(-1)^k\frac{A^{2k+1}}{(2k+1)!}+\cdots,
\]
\[
\cos(A)=I-\frac{A^2}{2!}+\frac{A^4}{4!}+\cdots+(-1)^k\frac{A^{2k}}{(2k)!}+\cdots,
\]
\[
\exp(A)=I+A+\frac{A^2}{2!}+\frac{A^3}{3!}+\cdots+\frac{A^k}{k!}+\cdots,
\]
其中 $I$ 是 $n\times n$ 单位矩阵。
\end{enumerate}

为应用于微分方程，引入变量 $t$。对 $A\in M_n(K)$，矩阵 $At$ 由 $A$ 中每个条目乘以 $t$ 得到。因此在收敛点上，$t\mapsto G(At)$ 是从 $K$ 的子集到 $M_n(K)$ 的函数。特别地，$\sin(At)$、$\cos(At)$ 与 $\exp(At)$ 对所有 $t\in K$ 收敛。

\item 设 $P$ 为非奇异 $n\times n$ 矩阵。
\begin{enumerate}[label=(\alph*)]
\item 证明
\[
PG(At)P^{-1}=G(PAtP^{-1})=G(PAP^{-1}t).
\]
这说明在相差一个基变换的意义下，只需对标准形矩阵计算 $G(At)$。提示：先对部分和取极限得到第一个等式；第二个等式由标量矩阵 $tI$ 与任意矩阵可交换立即得出。
\item 证明若 $A$ 是矩阵 $A_1,A_2,\ldots,A_m$ 的直和，则 $G(At)$ 是 $G(A_1t),G(A_2t),\ldots,G(A_mt)$ 的直和。
\item 证明若 $Z$ 是对角矩阵，对角元为 $z_1,z_2,
\ldots,z_n$，则 $G(Zt)$ 是对角矩阵，对角元为 $G(z_1t),G(z_2t),
\ldots,G(z_nt)$。
\end{enumerate}

矩阵 $\exp(A)$ 常记为 $e^A$。下面三题导出初等 Jordan 矩阵 $J$ 的 $\exp(Jt)$ 公式。

\item 证明若 $A$ 与 $B$ 可交换，则
\[
\exp(A+B)=\exp(A)\exp(B).
\]
提示：把 $A,B$ 看作可交换不定元，比较左边幂级数与右边两个幂级数的乘积。

\item 用上一题证明：若 $M$ 为任意矩阵，$\lambda\in K$，则
\[
\exp(\lambda It+M)=e^{\lambda t}\exp(M).
\]

\item 令 $N$ 为第一超对角线上全为 $1$、其余条目全为 $0$ 的 $r\times r$ 矩阵。计算如下幂零矩阵的指数：若
\[
Nt=\begin{pmatrix}
0&t&0&\cdots&0\\
0&0&t&\cdots&0\\
\vdots&\vdots&\ddots&\ddots&\vdots\\
0&0&\cdots&0&t\\
0&0&\cdots&0&0
\end{pmatrix},
\]
则
\[
\exp(Nt)=
\begin{pmatrix}
1&t&\frac{t^2}{2!}&\cdots&\frac{t^{r-1}}{(r-1)!}\\
0&1&t&\cdots&\frac{t^{r-2}}{(r-2)!}\\
\vdots&\vdots&\ddots&\ddots&\vdots\\
0&0&\cdots&1&t\\
0&0&\cdots&0&1
\end{pmatrix}.
\]
推出若 $J$ 是特征值为 $\lambda$ 的 $r\times r$ 初等 Jordan 矩阵，则
\[
\exp(Jt)=
\begin{pmatrix}
 e^{\lambda t}&te^{\lambda t}&\frac{t^2}{2!}e^{\lambda t}&\cdots&\frac{t^{r-1}}{(r-1)!}e^{\lambda t}\\
0&e^{\lambda t}&te^{\lambda t}&\cdots&\frac{t^{r-2}}{(r-2)!}e^{\lambda t}\\
\vdots&\vdots&\ddots&\ddots&\vdots\\
0&0&\cdots&e^{\lambda t}&te^{\lambda t}\\
0&0&\cdots&0&e^{\lambda t}
\end{pmatrix}.
\]
提示：由于 $Nt$ 幂零，$\exp(Nt)$ 是 $Nt$ 的多项式。计算 $\exp(Jt)$ 时写 $Jt=\lambda It+Nt$，再用练习 44。

若 $A\in M_n(K)$，$P$ 是基变换矩阵，使 $P^{-1}AP$ 为 Jordan 标准形。设 $P^{-1}AP$ 是初等 Jordan 矩阵 $J_1,
\ldots,J_m$ 的直和。由前面练习（取 $t=1$），先写 $E=\exp(P^{-1}AP)$ 为 $\exp(J_1),
\ldots,
\exp(J_m)$ 的直和，再变回原基，即得 $\exp(A)=PEP^{-1}$。

\item 对本节例 3 中给出的 $4\times4$ 矩阵 $D$ 与 $P$：
\[
D=\begin{pmatrix}
1&2&-4&4\\
2&-1&4&-8\\
1&0&1&-2\\
0&1&-2&3
\end{pmatrix},
\qquad
P=\begin{pmatrix}
0&1&2&0\\
2&0&-2&1\\
1&0&0&0\\
0&0&1&0
\end{pmatrix},
\]
证明
\[
E=\begin{pmatrix}
e&e&0&0\\
0&e&0&0\\
0&0&e&e\\
0&0&0&e
\end{pmatrix},
\qquad
\exp(D)=\begin{pmatrix}
e&2e&-4e&4e\\
2e&-e&4e&-8e\\
e&0&e&-2e\\
0&e&-2e&3e
\end{pmatrix}.
\]

\item 计算下列每个矩阵的指数：
\begin{enumerate}[label=(\alph*)]
\item 本节例 2 中的矩阵 $A$；
\item 练习 4 中的矩阵（在那里已经求了 Jordan 标准形和一个基变换矩阵）；
\item 练习 16 中的矩阵。
\end{enumerate}

\item 证明 $\exp(0)=I$（其中 $0$ 为零矩阵，$I$ 为单位矩阵）。推出对所有 $A\in M_n(K)$，$\exp(A)$ 非奇异，且逆矩阵为 $\exp(-A)$。

\item 证明
\[
\det(\exp(A))=e^{\operatorname{tr}(A)},
\]
其中 $\operatorname{tr}(A)$ 是 $A$ 的迹。

\item 固定 $A\in M_n(K)$。证明映射
\[
K\longrightarrow GL_n(K),
\qquad
 t\longmapsto \exp(At)
\]
是群同态（这里 $K$ 是域的加法群）。注意这推广了熟悉的指数映射 $K\to K^\times$，即 $n=1$ 情形。子群 $\{\exp(At)\mid t\in K\}$ 称为 $GL_n(K)$ 的一个\emph{单参数子群}。这些子群和指数映射在 Lie 群理论中很重要。

设 $G(x)$ 是收敛半径为无穷大的幂级数，并固定矩阵 $A\in M_n(K)$。矩阵 $G(At)$ 的条目是变量 $t$ 的函数。令 $c_{ij}(t)$ 是 $G(At)$ 的 $(i,j)$ 条目。$G(At)$ 对 $t$ 的导数记为 $\frac{d}{dt}G(At)$，定义为各条目逐项求导所得矩阵。

\item 建立下列导数性质：
\begin{enumerate}[label=(\alph*)]
\item 若 $G(x)=\sum_{k=0}^\infty \alpha_kx^k$，则
\[
\frac{d}{dt}G(At)=A\sum_{k=1}^\infty k\alpha_k(At)^{k-1}.
\]
\item 若 $v$ 是条目为常数的 $n\times1$ 矩阵，则
\[
\frac{d}{dt}\bigl(G(At)v\bigr)=\left(\frac{d}{dt}G(At)\right)v.
\]
\end{enumerate}

\item 由上一题 (a) 推出
\[
\frac{d}{dt}\exp(At)=A\exp(At).
\]

现在令 $y_1(t),
\ldots,y_n(t)$ 为实变量 $t$ 的可微函数，它们满足以下常系数一阶线性微分方程组：
\[
\begin{aligned}
y_1'&=a_{11}y_1+a_{12}y_2+\cdots+a_{1n}y_n,\\
y_2'&=a_{21}y_1+a_{22}y_2+\cdots+a_{2n}y_n,\\
&\ \vdots\\
y_n'&=a_{n1}y_1+a_{n2}y_2+\cdots+a_{nn}y_n.
\end{aligned}\tag{$*$}
\]
令 $A=(a_{ij})$，则可写为 $y'=Ay$。一个条目为 $t$ 的函数、其列向量为方程组 $(*)$ 的线性无关解的 $n\times n$ 矩阵，称为 $(*)$ 的\emph{基本矩阵}。由微分方程理论，$(*)$ 的解向量空间为 $n$ 维，其基本矩阵的列构成全部解空间的一组基。

\item 证明 $\exp(At)$ 是 $(*)$ 的基本矩阵。并证明：若 $C$ 是 $n\times1$ 常数列向量，其条目为 $y_1(0),
\ldots,y_n(0)$，则
\[
y(t)=\exp(At)C
\]
是满足初值条件 $y(0)=C$ 的 $(*)$ 的特解。注意这推广了 $1$ 维方程 $y'=ay$ 的结果：$e^{at}$ 是解空间的一组基，且满足 $y(0)=c$ 的唯一解是 $y=ce^{at}$。提示：使用前面练习。

\item 证明若 $M$ 是 $(*)$ 的基本矩阵，$Q\in M_n(K)$ 非奇异，则 $MQ$ 也是 $(*)$ 的基本矩阵。提示：$MQ$ 的列是 $M$ 的列的线性组合。

现在用前两题解决具体微分方程组：给定 $(*)$ 中的矩阵 $A$，求基变换矩阵 $P$ 使 $B=P^{-1}AP$ 为 Jordan 标准形。则 $\exp(At)=P\exp(Bt)P^{-1}$ 是 $(*)$ 的基本矩阵；由上一题，$P\exp(Bt)$ 也是基本矩阵。这样可用练习 45 后的方法计算 $\exp(Bt)$，且不必求 $P^{-1}$。

例如，若 $A=D$ 且 $P$ 为练习 46 中矩阵，则 $B=P^{-1}AP$ 由两个特征值为 $1$ 的 $2\times2$ Jordan 块组成。因此 $y'=Ay$ 的一个基本矩阵为
\[
P\exp(Bt)=
\begin{pmatrix}
0&1&2&0\\
2&0&-2&1\\
1&0&0&0\\
0&0&1&0
\end{pmatrix}
\begin{pmatrix}
e^t&te^t&0&0\\
0&e^t&0&0\\
0&0&e^t&te^t\\
0&0&0&e^t
\end{pmatrix}
=
\begin{pmatrix}
0&e^t&2e^t&2te^t\\
2e^t&2te^t&-2e^t&e^t(1-2t)\\
e^t&te^t&0&0\\
0&0&e^t&te^t
\end{pmatrix}.
\]
由此得方程组
\[
\begin{aligned}
y_1'&=y_1+2y_2-4y_3+4y_4,\\
y_2'&=2y_1-y_2+4y_3-8y_4,\\
y_3'&=y_1+y_3-2y_4,\\
y_4'&=y_2-2y_3+3y_4
\end{aligned}
\]
的一般解
\[
\begin{pmatrix}y_1\\y_2\\y_3\\y_4\end{pmatrix}
=\alpha_1\begin{pmatrix}0\\2e^t\\e^t\\0\end{pmatrix}
+\alpha_2\begin{pmatrix}e^t\\2te^t\\te^t\\0\end{pmatrix}
+\alpha_3\begin{pmatrix}2e^t\\-2e^t\\0\\e^t\end{pmatrix}
+\alpha_4\begin{pmatrix}2te^t\\e^t(1-2t)\\0\\te^t\end{pmatrix},
\]
其中 $\alpha_1,
\ldots,
\alpha_4\in K$ 任意。

\item 在 (a)--(c) 各部分中，对给定系数矩阵 $A$ 的方程组 $(*)$ 求一个基本矩阵。
\begin{enumerate}[label=(\alph*)]
\item $A$ 为练习 47(a) 中的矩阵。
\item $A$ 为练习 47(b) 中的矩阵。
\item $A$ 为练习 47(c) 中的矩阵。
\end{enumerate}

\item 考虑方程组 $(*)$，其系数矩阵 $A$ 为练习 46 中列出的矩阵 $D$，其基本矩阵已在上一题前算出。求满足初始条件 $y_i(0)=1$（$i=1,2,3,4$）的 $(*)$ 的特解。

接下来研究 $(*)$ 的一个特殊情形。给定常系数 $n$ 阶线性微分方程
\[
y^{(n)}+a_{n-1}y^{(n-1)}+\cdots+a_1y'+a_0y=0,\tag{$**$}
\]
其中 $y^{(k)}$ 表示 $y$ 的第 $k$ 阶导数，$y^{(0)}=y$。令 $y_i=y^{(i-1)}$，$1\leq i\leq n$，可把它化为一阶线性方程组。$(**)$ 的解空间的一组基可由该线性系统的基本矩阵得到：基本矩阵每一列的第 $(1,1)$ 项给出 $(**)$ 的一个解，第一行的 $n$ 个函数组成 $(**)$ 的解空间基。

\item 证明由 $(**)$ 得到的一阶系统的系数矩阵 $A$ 是多项式
\[
x^n+a_{n-1}x^{n-1}+\cdots+a_1x+a_0
\]
的伴随矩阵的转置。

\item 用上述方法求下列微分方程的解空间基：
\begin{enumerate}[label=(\alph*)]
\item $y'''-3y'+2y=0$；
\item $y''''+4y'''+6y''+4y'+y=0$。
\end{enumerate}

微分方程组
\[
\begin{aligned}
y_1'&=F_1(y_1,y_2,
\ldots,y_n),\\
y_2'&=F_2(y_1,y_2,
\ldots,y_n),\\
&\ \vdots\\
y_n'&=F_n(y_1,y_2,
\ldots,y_n)
\end{aligned}
\]
称为\emph{自治系统}，简记为 $y'=F(y)$，其中 $F=(F_1,
\ldots,F_n)$。自治意味着右端不显含时间变量 $t$。常系数线性系统 $(*)$ 是其中每个 $F_i$ 都为线性函数的特殊情形。

常数解 $y=c$（其中 $c=(c_1,
\ldots,c_n)^t$）满足 $y'=F(y)$ 当且仅当
\[
0=F_i(c_1,
\ldots,c_n),
\qquad i=1,2,
\ldots,n.
\]
这些常数解称为\emph{稳态}。若每个解都趋于稳态 $c$，即对任意解 $y(t)$，都有 $\lim_{t\to\infty}y_i(t)=c_i$，则称该稳态\emph{全局渐近稳定}。在线性自治系统中，解空间为向量空间，唯一常数解是零解。下一题给出零解全局渐近稳定的一个充分条件。

\item 证明：若 $A$ 的所有特征值实部均为负，则 $(*)$ 的零解 $y_i(t)=0$（$i=1,
\ldots,n$）全局渐近稳定。若不熟悉复指数函数的行为，可假设所有特征值均为实数（因此为负实数）。利用解的显式形式说明它们均在 $t\to\infty$ 时趋于零。
\end{enumerate}
